% lmcs, 2018.5.2, 2019.9.9, 2020.2.10
\documentclass{lmcs}
%\usepackage{fouriernc}
%\setmathfont[FakeBold=0.5]{Latin Modern Math}

\usepackage{wasysym}
\usepackage{hyperref}
\usepackage{graphicx}
\usepackage{amssymb,amsmath,graphicx,color}
\usepackage{mymath,proof,latexsym}
\usepackage{xcolor}
\usepackage[pdftex,outline]{contour}
\usepackage{tikz}
%%%%%%%%%%%%%%%%%%%%%%%%%%%%%%%%%%%%
% WE ADD THE COMMANDS 
% theorem lemma proposition corollary definition  
%%%%%%%%%%%%%%%%%%%%%%%%%%%%%%%%%%%%
\newtheorem{theorem}{Theorem}[section]
\newtheorem{lemma}[theorem]{Lemma}
\newtheorem{proposition}[theorem]{Proposition}
\newtheorem{corollary}[theorem]{Corollary}
\newtheorem{definition}[theorem]{Definition}

%\newenvironment{proof}[1][Proof]{\begin{trivlist}
%\item[\hskip \labelsep {\bfseries #1}]}{\end{trivlist}}

\newenvironment{example}[1][Example]
{\begin{trivlist} \item[\hskip \labelsep {\bfseries #1}]}{\end{trivlist}}

\newenvironment{remark}[1][Remark]
{\begin{trivlist}\item[\hskip \labelsep {\bfseries #1}]}{\end{trivlist}}


\begin{document}
\sloppy 
\hbadness=10000
\vbadness=10000

\def\cal{\mathcal} % for LMCS
\def\calH{{\cal H}}

% journal, pa2

\def\PAtwoID{{\rm PA}_2{\rm ID}}
\def\CPAtwoID{{\rm CPA}_2{\rm ID}}
\def\AxiomPA{{\rm AxiomPA}}

% cmcs

\def\Co{{\rm co}}
\def\Axiom{{\rm Axiom}}
\def\Wk{{\rm Wk}}
\def\Cut{{\rm Cut}}
\def\BS{{\rm BS}}
\def\Bit{{\rm Bit}}
\def\LK{{\rm LK}}
\def\Head{{\rm head}}
\def\Tail{{\rm tail}}

% equivalence cyclic proofs

%\def\Uor{\displaystyle \mathop {\widetilde\bigvee}}
\def\Uor{\biguplus}
\def\CLKIDomega{{{\bf CLKID}^\omega}}
\def\CLKIDPA{{{\bf CLKID}^\omega+{\bf PA}}}
\def\LKIDPA{{{\bf LKID}+{\bf PA}}}
\def\PA{{\bf PA}}

% equiv note

\def\Choose{{\rm Choose}}
\def\Asym{{\rm Asym}}
\def\Peano{{\bf Peano}}
\def\Code#1{{\lceil{#1}\rceil}}
\def\Ineq{{\rm Ineq}}
\def\MonoPath{{\rm MonoPath}}
\def\MonoSeq{{\rm MonoSeq}}
\def\BoundedPath{{\rm BoundedPath}}
\def\InfPath{{\rm InfPath}}
\def\InfMonoPath{{\rm InfMonoPath}}
\def\InfMonoSeq{{\rm InfMonoSeq}}
\def\HomSeq{{\rm HomSeq}}
\def\InfHomSeq{{\rm InfHomSeq}}

\def\ErdosTree{{\rm ErdosTree}}
\def\KB{{\rm KB}}
\def\MonoList{{\rm MonoList}}
\def\Seq{{\rm Seq}}
\def\Univ{{\rm univ}}
\def\Infinite{{\rm Infinite}}
\def\S{{\bf S}}
\def\N{{\bf N}}
\def\Ind{{\rm Ind}}
\def\LKID{{\bf LKID}}
\def\PA{{\bf PA}}
\def\LKIDExt{{\bf LKIDExt}}
\def\CLKID{{\bf CLKID}}

% LICS

\def\Over#1#2{\deduce{#2}{#1}}
\def\List{{\rm List}}
\def\ListX{{\rm ListX}}
\def\ListO{{\rm ListO}}
\def\ListE{{\rm ListE}}
\def\Nl{{\rm nl}}

% CSLID

\def\SLLID{{\rm SLID}}
\def\SLID{{\rm SLID}}
\def\CSLID{{\rm CSLID}}
\def\CSLIDomega{{{\rm CSLID}\omega}}
\def\Eclass{{\rm Eclass}}
\def\Eq{{\rm Eq}}
\def\Deq{{\rm Deq}}
\def\Satom{{\rm Satom}}
\def\Cells{{\rm Cells}}
\def\Roots{{\rm Roots}}
\def\Elim{{\rm Elim}}
% \def\Case#1{{\rm Case-$#1$}}
\def\Case{{\rm Case}}
\def\Unfold{{\rm Unfold}}
\def\Pred{{\rm Pred}}
\def\Start{{\rm Start}}
\def\Unsat{{\rm Unsat}}
\def\Split{{\rm Split}}
\def\Underscore{\underline{\phantom{x}}}
\def\Rootcell{{\rm rootcell}}
\def\Rootshape{{\rm rootshape}}
\def\Jointleaf{{\rm jointleaf}}
\def\Jointleafimplicit{{\rm jointleafimplicit}}
\def\Jointnode{{\rm jointnode}}
\def\Directjoint{{\rm directjoint}}

% cyclicnote2

\def\Range{{\rm Range}}
\def\Xstart{X_{{\rm start}}}
\def\kstart{k_{{\rm start}}}
\def\Leaves{{\rm Leaves}}
\def\PureDist{{\rm PureDist}}
\def\Pure{{\rm Pure}}
\def\Identity{{\rm Identity}}
\def\Subst{{\rm Subst}}

% weak establish

\def\Connected{{\rm Connected}}
\def\Eststablished{{\rm Eststablished}}
\def\Valued{{\rm Valued}}

% monadic new

\def\SLRDbtw{{\rm SLRD}_{btw}}
\def\BTW{{{\rm BTW}}}
\def\Loc{{{\rm Loc}}}
\def\Ne{{\ne}}
\def\Nil{{{\rm nil}}}
\def\eqDef{=_{{\rm def}}}
\def\Inf#1{{\infty_{#1}}}
\def\Stores{{\rm Stores}}
\def\SVars{{{\rm SVars}}}
%\def\Val{{{\rm Val}}}
\def\MSO{{{\rm MSO}}}
\def\Sep{{{\rm Sep}}}
\def\Septwo{{{\rm SLMI}}}
\def\SLMI{{{\rm SLMI}}}
\def\Sepinf{{\rm Sep}\infty}
\def\THeaps{{\rm THeaps}}
\def\Root{{\rm Root}}
\def\TG{{\rm TGraph}}

% monadic

\def\nil{{{\rm nil}}}
\def\eqDef{=_{{\rm def}}}
\def\Inf#1{{\infty_{#1}}}
\def\Stores{{\rm Stores}}
\def\SVars{{{\rm SVars}}}
\def\Val{{{\mathcal V}}}
\def\MSO{{{\rm MSO}}}
\def\Sep{{{\rm sep}}}
\def\THeaps{{\rm THeaps}}
\def\Root{{\rm Root}}
\def\TG{{\rm TGraph}}

\def\Tilde{\widetilde}
\def\Bar{\overline}
\def\Lequiv{\Longleftrightarrow}
\def\Lto{\Longrightarrow}
\def\Lfrom{\Longleftarrow}
\def\Norm{{\rm Norm}}
\def\Noshare{{\rm Noshare}}
\def\Roots{{\rm Roots}}
\def\Forest{{\rm Forest}}
\def\Var{{\rm Var}}
\def\Range{{\rm Range}}
\def\Cell{{\rm Cell}}
\def\Tree{{\rm Tree}}
\def\Switch{{\rm Switch}}
\def\All{{\rm All}}
\def\To{\leadsto}
\def\tree{{\rm tree}}
\def\Paths{{\rm Paths}}
\def\Finite{{\rm Finite}}
\def\Const{{\rm Const}}
\def\Leaf{{\rm Leaf}}
\def\LeastElem{{\rm LeastElem}}
\def\LeastIndex{{\rm LeastIndex}}
\def\WSnS{{{\rm WSnS}}}
\def\Expand{{{\rm Expand}}}

% previous

\long\def\J#1{} % skip
\def\T#1{\hbox{\color{green}{$\clubsuit #1$}}}
\def\W#1{\hbox{\color{Orange}{$\spadesuit #1$}}}

\def\Node{{{\rm Node}}}
\def\LL{{{\rm LL}}}
\def\DSN{{{\rm DSN}}}
\def\DCL{{{\rm DCL}}}
\def\LS{{{\rm LS}}}
\def\Ls{{{\rm ls}}}

\def\FPV{{{\rm FPV}}}
\def\Lfp{{{\rm lfp}}}
\def\IsHeap{{{\rm IsHeap}}}

\def\Equiv{\quad \equiv\quad }
\def\Null{{{\rm null}}}
\def\Emp{{{\rm emp}}}
\def\If{{{\rm if\ }}}
\def\Then{{{\rm \ then\ }}}
\def\Else{{{\rm \ else\ }}}
\def\While{{{\rm while\ }}}
\def\Do{{{\rm \ do\ }}}
\def\Cons{{{\rm cons}}}
\def\Dispose{{{\rm dispose}}}
\def\Vars{{{\rm Vars}}}
\def\Locs{{{\rm Locs}}}
\def\States{{{\rm States}}}
\def\Heaps{{{\rm Heaps}}}
\def\FV{{{\rm FV}}}
\def\True{{{\rm true}}}
\def\False{{{\rm false}}}
\def\Dom{{{\rm Dom}}}
\def\Abort{{{\rm abort}}}
\def\New{{{\rm New}}}
\def\W{{{\rm W}}}
\def\Pair{{{\rm Pair}}}
\def\Lh{{{\rm Lh}}}
\def\lh{{{\rm lh}}}
\def\Elem{{{\rm Elem}}}
\def\EEval{{{\rm EEval}}}
% \def\BEval{{{\rm BEval}}} % not used 
\def\PEval{{{\rm PEval}}}
\def\HEval{{{\rm HEval}}}
\def\EVal{{{\rm Eval}}}
\def\Domain{{{\rm Domain}}}
\def\Exec{{{\rm Exec}}}
\def\Store{{{\rm Store}}}
\def\Heap{{{\rm Heap}}}
\def\Storecode{{{\rm Storecode}}}
\def\Heapcode{{{\rm Heapcode}}}
\def\Lesslh{{{\rm Lesslh}}}
\def\Addseq{{{\rm Addseq}}}
\def\Separate{{{\rm Separate}}}
\def\Result{{{\rm Result}}}
\def\Lookup{{{\rm Lookup}}}
\def\ChangeStore{{{\rm ChangeStore}}}
\def\ChangeHeap{{{\rm ChangeHeap}}}
\def\Wand{\mathbin{\hbox{\hbox{---}$*$}}}
\def\Eval#1{[\kern -1pt[{#1}]\kern -1pt]}
\def\Vec{\overrightarrow}
%\def\Vec#1{{\bf #1}}

\def\Tilde{\widetilde}
\def\Break{\hfil\break\hbox{}}

\newcommand{\freshcopy}           [1]                      {  {\tilde{#1}} }
%\newcommand{\freshcopy}           [1]                      {  \displaystyle{\mathop{#1}^{\sim}} }
\newcommand{\bfColor} [2]     {{\bf \color{#1}{#2}}}
\newcommand{\arity}               {{\tt ar}}          
\newcommand{\trace}               {{\tt tr}} 
\newcommand{\stat}                 {{\tt stat}} 
\newcommand{\Stat}                {{\tt S}}     
\newcommand{\prog}               {{\tt p}} 
\newcommand{\Prog}               {{\tt P}}   
\newcommand{\fun}                 {{\tt fun}}          
\newcommand{\quotationMarks} [1]  {{\textquotedblleft #1\textquotedblright}}
\newcommand{\universe} [1]  {{| #1 |}}
\newcommand{\Rule}               {{\tt Rule}} 
\newcommand{\Path}               {{\tt Path}} 
\newcommand{\pathBetween} {{\tt path}} 
\newcommand{\Trace}               {{\tt Trace}} 
\newcommand{\Cyclic}              {{\tt Cyclic}}
\newcommand{\Cycle}              {{\tt Cycle}}
\newcommand{\dom}              {{\tt dom}}
\newcommand{\codom}              {{\tt codom}}
\newcommand{\Periodic}              {{\tt Periodic}}
\newcommand{\setR}              {{\mathcal R}}
\newcommand{\setL}              {{\mathcal L}}
\newcommand{\setG}              {{\mathcal G}}
\newcommand{\Pfin}      {{ {\mathcal R}_{\mbox{\tiny fin}} }}
\newcommand{\Card}              {{\tt Card}}
\newcommand{\restr}               {{\lceil}}
\newcommand{\comp}              {{\circ}}
\newcommand{\val}                   {{\tt val}}
\newcommand{\setP}                   {{\mathcal P}}
\newcommand{\Rprog}               {{\star}}
\newcommand{\RProg}               {{\bigstar}}
\newcommand{\length}               {{\mbox{\tt len}}}
\newcommand{\lex}                    {{\mbox{\tiny lex}}}
\newcommand{\setH}                  {{\mathcal H}}
\newcommand{\setS}                  {{\mathcal S}}
\newcommand{\setT}                  {{\mathcal T}}
\newcommand{\id}                      {{\tt id}}
\newcommand{\BrowerKleene}  {{\mbox{\tiny BK}}}
\newcommand{\conc}                 {{*}}




% MORE MACRO
\newcommand{\size} 		                    {{\tt s}}
\newcommand{\state} 		                 {{\tt S}}
\newcommand{\minStable}               {{\tt m}}
\newcommand{\unfold}                     {{\tt u}}
\newcommand{\height}                      {{\tt h}}
\newcommand{\setF} 		                     {{\mathcal F}}
\newcommand{\setU} 	                     {{\mathcal U}}
\newcommand{\cons}                         {{\tt cons}}
\newcommand{\newton} [2]               {{\genfrac(){0pt}{1} {#1}{#2}}}
%\newcommand{\quotationMarks} [1] {{‘‘#1’’}}
\newcommand{\MaxInd}                  {{\tt M}}
\newcommand{\Raising}                {{\tt Raising}}
\newcommand{\trans}                       {{\tt trans}}
\newcommand{\DomCard}               {{\tt DomCard}}
\newcommand{\grouping}                {{\succeq}}
\newcommand{\cod}                         {{\tt cod}}
%END MACRO

\title{A Poly-exponential Translation from Injective Cyclic to Inductive Proofs}
\date{}

\author[S. Berardi]{Stefano Berardi}
\address{Universit\`{a} di Torino,
Torino, Italy}
\email{stefano.berardi@unito.it}

%% mandatory lists of keywords 
\keywords{
proof theory
inductive definitions
cyclic proofs
infinite Ramsey theorem
}

%%%%%%%%%%%%%%%%%%%%%%%%%%%%%%%
% Titolo originario: An Intuitionistic and First Order proof that GTC implies Correctnes
%%%%%%%%%%%%%%%%%%%%%%%%%%%%%%%
% questo lavoro mio in data 16/07/2025 è diventato un lavoro su injectivity e prove
% cicliche, per CSL 2026.  Ci sono ora tre nomi, io, Osiero e de' Liguoro. Viene rifiutato
% in data 15 ottobre 2026. Un recensore, forse Amram, è stato particolarmente
% critico, sostiene che il lavoro non contiene alcun progresso rispetto al suo
% lavoro SCT con Niel Jones. In data 22 Ottobre provo a riprendere il lavoro del 
% 16/072025 per vedere se si può trasformare in un lavoro di proof-theory e 
% trasformazione di prove cicliche in prove ordinarie, che non possa venir recensito
% da amram.
%%%%%%%%%%%%%%%%%%%%%%%%%%%%%%%


\maketitle


\emph{Proof started in Torino, January 20, 2026. 
Ugo suggested adding complexity results to motivate the new algorithm}.

Mirai \cite{Mirai} provided a translation from from cyclic to inductive proofs. The proof builds
on an existing algorithm for generating ranking functions and 
requiring triple exponential time. A first improvement would be to switch to an equivalent
formulation of cyclic proofs with injective bicolored relation (i.e. fan-in free, with at most one
arrow entering in each node). 
In this case, the same algorithm provides a double exponential time \cite{Lee2} 
and an inductive proof in output of size $2^{O(n \log(n))}$.
In this paper we further improve the translation to a single exponential time
(poly-exponential time) , by introducing a new algorithm. The price is having an inductive proof 
in output of poly-exponential size, to be precise $2^{O(n^3)}$ size, instead than an ouput 
of size $2^{O(n \log(n))}$.

The translation starts from a cyclic proof $P$ with set $\setR$ of bicolored relations on a set 
$[1,\MaxInd]$. $\MaxInd$ is the maximum number of inductive formula in the left-hand side 
of some sequent in $\Pi$. 


\section{Quick Overview of the GTC, Injectivity and PDC Conditions}

\subsection{Bicolored Binary relations on a Finite Set}
A bicolored relation is any pair $(\Stat,\Prog)$ of binary relations on
$[1,\MaxInd]$ such that $\Stat \cap \Prog = \emptyset$. We say that any pair $(x,y) \in \Stat$
has the stationary color, and any pair $(x,y) \in \Prog$
has the progressing color. If $R= (\Stat,\Prog)$ is a bicolored relation, we define 
$\universe{R} = \Stat \cup \Prog$ the uncolored associated relation. 
We have a notion of composition for bicolored relations 
$(\Stat,\Prog)(\Stat',\Prog') = (\Stat'', \Prog'')$ with $\Prog'' = \Stat \Prog' \cup 
\Stat' \Prog \cup \Prog \Prog'$ and $\Stat'' = \Stat \Stat' \setminus \Prog''$. In the
case of colored relations associated to injective relations the definition of $\Stat''$
can be simplified to $\Stat'' = \Stat \Stat'$. Uncoloring and composition commute:
we have $\universe{RS} = \universe{R}\universe{S}$ for all colored relations $R$, $S$.

Each (possibly colored) relation $R$ is associated to an arrow $a \rightarrow b$ 
between nodes of $\Pi$, and any two relations $R:a \rightarrow b$, $S:b \rightarrow c$
are said to be composable, with composition $RS: a \rightarrow c$. 
We fix a $\setR$ of colored relations $R:a \rightarrow b$
in the paper, for some $R, a, b$. 
 
From now on we write $\N$ for the set of non-negative integers and $\N^+$ for the set of positive integers. For any $a,b \in \N$
we write $[a,b]$ for the segment $\{c \in \N | a \le c \le b\}$. A sub-segment $[a',b']$ of
$[a,b]$ is any $[a',b'] \subseteq [a,b]$, namely with $[a',b']=\emptyset$ or $a \le a', \le b' \le b$.

We consider binary relations on a finite set: for sake of simplicity we only 
consider the finite sets $[1,\MaxInd]$, for some $\MaxInd \in \N$ which we assume to be
fixed. 

We will abbreviate \quotationMarks{binary relations on $[1,\MaxInd]$} with
\quotationMarks{relations}. We write $\dom(R) = \{x | \exists y (x R y)\}$ and
$R(x)$ for the set $\{y | x R y\}$ of $R$-images of $x \in \dom(R)$. For colored
relation we define $\dom(R) = \dom(\universe{R})$. We call $\DomCard(R) = \Card(\dom(R))$ 
the domain-cardinality of $R$. We extend $\dom$ and $\DomCard$ to relation sequences by 
taking the product of the sequence: $\dom(R_1, \ldots, R_n) =  \dom(R_1 \ldots R_n)$ and 
$\DomCard(R_1, \ldots, R_n) = \DomCard(R_1 \ldots R_n)$.


\subsection{GTC, Injectivity and PDC Conditions}

\begin{definition}[Injectivity, PDC]
\label{definition-inj-PDC}
Assume $\setR^+$ is the closure of a family $\setR$ under colored composition, 
$R=(\Stat,\Prog)$ is a bi-colored relation on $[1,\MaxInd]$, and 
$X,Y \subseteq [1,\MaxInd]$ are sets of positive integer.

\begin{enumerate}
\item
\label{definition-inj-PDC-01}
$R$ is progressing from $X$ to $Y$ if and only if $\Prog \cap (X \times Y) \not = \emptyset$.

\item
\label{definition-inj-PDC-02}
The Progressing Domain Condition (PDC) for a bicolored relation $R : a \rightarrow a$ says: 
$R$ is progressing from $\dom(R)$ to $\dom(R)$.

\item
\label{definition-inj-PDC-03}
We say that closure $\setR^+$ of $\setR$
satisfies PDC if and only if for all $R: a \rightarrow a \in \setR$, $R$ satisfies PDC. 
\end{enumerate}
\end{definition}


\begin{definition}[The Global Trace Condition GTC]
\label{definition-GTC}
\begin{enumerate}
\item
\label{definition-GTC-01}
An infinite path $\pi$ in $\setR^+ $ is any infinite composable sequence
$$\{(\Stat_i, \Prog_i):  \alpha_i \rightarrow \alpha_{i+1}\}_{i \in \N} \subseteq \setR^+$$
\begin{enumerate}
\item
An infinite trace $\sigma$ in $\pi$ is any map $\sigma: [i_0, \infty[ \rightarrow [1,\MaxInd]$, 
for some $i_0 \in \N$, such that $\sigma(i) \Stat_i \cup \Prog_i \sigma(i+1)$ 
for all $i \in [i_0, \infty[$. 

\item
$\sigma$ is infinitely progressing in $\pi$ if and only if $\sigma(i) \Prog_i \sigma(i+1)$ 
for infinitely many $i \ge i_0$.
\end{enumerate}

\item
\label{definition-GTC-02}
We say that $\setR^+ $ satisfies GTC if and only if for all infinite composable sequences 
$\pi$ there is some infinitely progressing path $\sigma$ in $\pi$.

\end{enumerate}
\end{definition}

Remark that the Global Trace Condition (GTC) is a property of infinite composable
chains of $\setR^+$, while PDC is a property of single relations $R:a \rightarrow a \in \setR^+$,
which is extended to a property of  $\setR^+$. 

Through the paper assume that:

\begin{enumerate}
\item  
$\setR^+$ is a set of injective relations.
\item
$\setR^+$ satisfies PDC 
\end{enumerate}

PDC for injective relations is equivalent to GTC, while in general GTC is a stronger condition
than PDC. 
\\

\begin{example}
\label{example-PDC-GTC}
Let $R = (\{(1,1)\},\{(2,1)\})$ with $R$ stationary on $(1,1)$
and progressing on $(2,1)$. Let $\setR=\{R: a \rightarrow a\}$.
Then $R$ satisfies PDC %on the graph $\{(a,a)\}$, 
and $RR=R$, therefore  $\setR^+=\{R: a \rightarrow a\}$ and $\setR^+$ satisfies PDC. 
However, the only infinite traces in the infinite path 
$R: a \rightarrow a, R: a \rightarrow a, R: a \rightarrow a, \ldots$ are $(1,1,1,\ldots)$ and $(2,1,1,1,\ldots)$. The first trace never progresses, the second trace progresses once, 
on the first arrow $(2,1)$. 

So $\setR^+$ does \emph{not} satisfy GTC. The fact that GTC is false for $\setR^+$ 
is not in contradiction with what we said. Indeed, $R \in \setR^+$ is  \emph{not} injective,
given that $1$ has two $R$-counter-images, namely $1$ and $2$.
\end{example}

\begin{center}
\begin{tikzpicture}
\draw[->,blue, thick] (0,0) -- (1,0);
\draw[->,red, thick]   (0,1) -- (1,0);

\draw[->,red, thick]   (1.1,1) -- (2.1,0);
\draw[->,blue, thick] (1.1,0) -- (2.1,0);

\draw[->,red, thick]   (2.2,1) -- (3.2,0);
\draw[->,blue, thick] (2.2,0) -- (3.2,0);

\draw[->,red, thick] (3.3,1) -- (4.3,0);
\draw[->,blue, thick] (3.3,0) -- (4.3,0);

\filldraw[black] (0,0) circle (1pt) node[anchor=west]{1};
\filldraw[black] (0,1) circle (1pt) node[anchor=west]{2};

\filldraw[black] (1.1,0) circle (1pt) node[anchor=west]{1};
\filldraw[black] (1.1,1) circle (1pt) node[anchor=west]{2};

\filldraw[black] (2.2,0) circle (1pt) node[anchor=west]{1};
\filldraw[black] (2.2,1) circle (1pt) node[anchor=west]{2};

\filldraw[black] (3.3,0) circle (1pt) node[anchor=west]{1};
\filldraw[black] (3.3,1) circle (1pt) node[anchor=west]{2};


\filldraw[black] (4.4,0) circle (1pt) node[anchor=west]{1};
\filldraw[black] (4.4,1) circle (1pt) node[anchor=west]{2};

\filldraw[black] (5.5,0) circle (1pt) node[anchor=west]{$\ldots$};
\filldraw[black] (5.5,1) circle (1pt) node[anchor=west]{$\ldots$};

\node[] at (0,1.5) {$R$};
\node[] at (1.1,1.5) {$R$};
\node[] at (2.2,1.5) {$R$};
\node[] at (3.3,1.5) {$R$};
\node[] at (4.4,1.5) {$\ldots$};
\node[] at (5.5,1.5) {$\ldots$};

\end{tikzpicture}
\\
\emph{\small The only infinite traces are $(1,1,1,\ldots)$ and $(2,1,1,1,\ldots)$}
\end{center}

%Suppose we have a path $\pi$ from a companion $c$ to itself, moving back  through $c$
%for $k \ge 1$ times: we have $c \rightarrow c \rightarrow \ldots \rightarrow c$, with $k$ arrows.
%For any $i \in \{1, \ldots, k\}$, suppose the step $c \rightarrow c$ having number $i$
%is associated to some bicolored relation $R_i : c \rightarrow c$, $R_i \in \setR^+$. 
%The path $\pi$ itself is associated to the sequence $\vec{R} = R_1 \ldots, R_k$ of such 
%relations. For the translation, we have to associate a measure to each such sequence $\vec{R}$
%which is decreasing by extension.


\section{Infinite Traces of an infinite Path and Injectivity}
We first provide a picture of the traces of the infinite traces of an infinite path in the case of 
injective relations
$$
\pi = 
R_1:\alpha_1 \rightarrow \alpha_2, \ \ \ 
\ldots, \ \ \ 
R_k:\alpha_k \rightarrow \alpha_{k+1}, \ \ \ 
\ldots
$$
then we will sketch a decreasing measure on each finite prefix $R_1 \ldots, R_k$ of $\pi$. 
This section is informal and only intended as a plan for the paper.
\\ 

We assume having some map $\size:Nodes, [1,\MaxInd] \rightarrow \N^+$, assigning a positive
integer $\size(\alpha, x)$ to the inductive formula number $y$ in the node $\alpha$.
We assume that if $x$ progresses to $y$ when moving from node $\alpha$ to node $\beta$, 
then $\size(\alpha, x) > \size(\beta, x)$, while if $(x,y)$ is stationary then 
$\size(\alpha, x) = \size(\beta, x)$. 

When we study GTC, for each $i \in \N$ we only are interested to elements $x \in [1,\MaxInd]$
from which some infinite trace $x_i R_i x_{i+1} R_{i+1} x_{i+2} \ldots$ starts or
passes through. For each $i \in \N$, these elements are the elements of the infinite intersection
$$
I_i = \dom(R_i) \cap \dom(R_i R_{i+1}) \cap \dom(R_i R_{i+1} R_{i+2}) \ldots
$$
It is immediate to check that 
$
\dom(R_i) \supseteq 
\dom(R_i R_{i+1}) \supseteq 
\dom(R_i R_{i+1} R_{i+2}) \supseteq 
\ldots
$,
therefore $I_i = \bigcap_{j \ge i} \dom(R_i \ldots R_j)$. $I_i$ is an infinite intersection of a
weakly decreasing family of finite sets, therefore $I_i = \dom(R_i \ldots R_j)$ 
for some $j \ge i$. If we group the relations of $\pi$ in these groups
$R_i \ldots R_j$, we obtain an infinite path 
$$
\pi' = 
S_1:\beta_1 \rightarrow \beta_2, \ \ \ 
S_2:\beta_2 \rightarrow \beta_3, \ \ \  
S_3:\beta_3 \rightarrow \beta_4, \ \ \  
\ldots, \ \ \ 
S_m:\beta_m \rightarrow \beta_{m+1}, \ \ \ 
\ldots
$$ 
which we call \emph{stable}. This path has a simpler behaviour, it has the property
$$
\dom(S_i) = 
\dom(S_i S_{i+1}) =
\dom(S_i S_{i+1} S_{i+2}) =
\ldots
$$
Each $x \in \dom(S_i)$ is the first element of one or more infinite traces in the path $\pi'$,
to put otherwise, each $x \in \dom(S_i)$ has some $S_i$-image $y \in \dom(S_{i+1})$
Furthermore, in the case of \emph{injective} relations, these images are distinct, therefore
$\Card(\dom(S_i)) \le \Card(\dom(S_{i+1})$. Let $c_i = \Card(\dom(S_i))$:then 
$c_1 \le c_2 \le \ldots \le c_m \le \ldots$. 
For all $i \in \N$ we have $c_i = \Card(\dom(S_i)) \le \MaxInd$, therefore
we can decompose $\pi'$ in finitely many sub-sequences, the last one infinite, such that  
$$
c_{a_1} = \ldots = c_{b_1}   < 
c_{a_2} = \ldots = c_{b_2}   < 
\ldots                                       <
c_{a_h} = \ldots 
$$
We decompose $\pi' = S_1, S_2, \ldots$ in the same way:
$S_{a_1}, \ldots, S_{b_1}$, \ldots, $S_{a_{h-1}}, \ldots, S_{b_{h-1}}$, $S_{a_{h}}, \ldots$. 
We call these subsequences the \emph{components} of $\pi$.
Again by injectivity, we can prove that $S_{a_1}$ is a bijection from 
$\dom(S_{a_1})$ to $\dom(S_{a_1 + 1})$, $S_{a_1 + 1}$ is a bijection from 
$\dom(S_{a_1 + 1})$ to $\dom(S_{a_1 + 2})$ and \ldots and 
$S_{b_1 - 1}$ is a bijection from $\dom(S_{b_1 - 1})$ to $\dom(S_{b_1})$.
The same is true for all sub-sequences, including the infinite subsequence $S_{a_{h}}, \ldots$.

This means that eventually the traces all follow the bijections in $S_{a_{h}}, \ldots$,
and from the index $a_h$ on there are at most $c$ traces, with $c = c_{a_h}$
the common cardinality of all $\dom(S_j)$ for $j \ge a_{h}$. There are at most 
$\newton{\MaxInd}{c}$ different sets of
cardinality $c$, therefore for some $j > k \ge a_h$ we have $\dom(S_j) = \dom(S_k)$
and $S_j \ldots S_{k-1}$ is a bijection from $\dom(S_j)$ to $\dom(S_k) = \dom(S_j)$.
Using PDC, we will prove that $S_j \ldots S_{k-1}$ is progressing from some 
$x \in \dom(S_j)$ to some $y \in \dom(S_k)$. This progression happens at least once each 
$\newton{\MaxInd}{c}$ steps.

This first suggesta using an induction with last parameter 
$$
\Sigma_{x \in \dom(S_n)} \size(\beta_n,x)
$$ 
Whenever some $x \in \dom(S_n)$ progresses then the sum decreases. 
However, this idea is incomplete in several ways, as we will see.
\begin{enumerate}

\item
First, given a finite prefix $R_1, \ldots, R_n$ of $\pi$, we ignore whether the last component 
$S_{a_{h}}, \ldots, S_m$ of the stable version $S_1, \ldots, S_m$ of $R_1, \ldots, R_n$ 
is the last component of the infinite sequence $\pi$ or not. The grouping required to produce a 
stable relation sequence may erase one or more components to the end. 
We can fix this objection by taking nested induction on:
$$
\Sigma_{x \in \dom(S_{b_1})} \size(\beta_{b_1},x), \ \ \ 
\ldots, \ \ \ 
\Sigma_{x \in \dom(S_{b_h})} \size(\beta_{b_h},x)
$$
When a component is erased we use induction on the index for the last non-erased component
$\Sigma_{x \in \dom(S_{b_i})} \size(\beta_{b_i},x)$.

\item
Second, the last relation $S_m$ makes a special case.
If $x \in \dom(S_{b_h})$ has a set $S_{b_h}(x)$ of one or more $S_n$-images, 
then any of them can be the next element of the infinite progressing trace. 
The worst case from the viewpoint of the induction is when the next element is 
some $y \in S_{b_h}(x)$ such that $\size(\beta_{b_h},y)$ is maximum. 
This suggests switching to the indexes
$$
\Sigma_{x \in \dom(S_{b_i})} \max(\{\size(\beta_{b_i+1},y) | y \in S_{b_i}(x)\}
$$
expressing such worst-case possibility.

\item
Third, we can have up to $\newton{\MaxInd}{c_{a_i}}$ steps until the index number $i$
decreases (even if such a bound looks a vast exaggeration). We need an auxiliary index 
$t \in [0,\newton{\MaxInd}{c_{a_i}}]$ which, whenever the index number $i$ is the last index,
counts how many steps we can still have 
until the index number $i$ decreases. Such an index decreases by $1$ at each step
in which the index number $i$ is the last index and it does not decrease.

\end{enumerate}

In the rest of the paper we will realize such a plan. We first to derive some elementar properties 
of the sequences of uncolored injective relations on $[1,\MaxInd]$ whic are required to
prove that our index decrease. These properties are combinatorial properties and they 
have nothing to do with colors, nor with cyclic proofs.


%19:00 10/02/2026


\section{Domain-Stable Sequences and Grouping}


We define the grouping relation between sequences of (uncolored) relations. 
An example: a possible grouping of $\vec{R} = R_1, R_2, R_3, R_4, R_5, R_6$ 
is the following list $\vec{T} = T_1, T_2, T_3$ of $3$ products:
 
$$ 
R_1S_2, \ \ \ R_3, \ \ \ R_4 R_5 R_6
$$

\begin{definition}[The Grouping Relation]
Assume $\vec{T} = T_1, \ldots, T_m$ and 
$\vec{R} = R_1, \ldots, R_n$. We define the grouping relation $\vec{T} \grouping \vec{R}$ 
by induction on $n$ as follows.
\begin{enumerate}

\item
If $n=0$, then $\vec{T}$ = empty sequence = $\vec{R}$.

\item
Suppose $\vec{T} = T_1, \ldots, T_m \grouping \vec{R}$. 
Then 
\begin{enumerate}
\item $T_1, \ldots, T_{i-1}, (T_i \ldots T_m R) \ \grouping \ \vec{R},R$, for all $i \in [1,m]$
\item $T_1, \ldots, T_{m}, R \ \grouping \ \vec{R},R$ 
\end{enumerate}
\end{enumerate}
\end{definition}

To put otherwise, the groupings of $\vec{R} = R_1, \ldots, R_n$ are all and only 
sequences of non-empty products of consecutive elements of $\vec{R}$, of the form: 
$$
(R_{a_1} \ldots R_{b_1}), \ \ \ 
(R_{a_2} \ldots R_{b_2}), \ \ \ 
\ldots, \ \ \ 
(R_{a_{h}} \ldots R_{b_{h}})
$$
The segments $[a_1,b_1], \ldots, [a_{h},b_{h}]$ are a decomposition
of $[1,n]$ into a disjoint union for some increasing sequence 
$0 < b_1 < b_2 <  \ldots < b_h = n$. We necessarily have
$a_0=1$ and $a_{i+1} = b_i + 1$ for all $i \in [1,h-1]$. 

Some trivial examples. 
\begin{enumerate}
\item
We can have $h=0$, $b_1=n$ and the trivial grouping consisting of a single product  
$(R_1 \ldots R_n)$ of all relations. 
\item
We can have $h=m$, $b_1=1$, \ldots, $b_m=n$ and the trivial grouping consisting of the 
sequence $R_1, \ldots, R_n$ itself, considered a sequence of unary products.
\end{enumerate}

We first prove some domain inclusions and equalities 
which are valid for all sequences of relations, then we isolate the \quotationMarks{stable}
grouping, in which more domain inclusions are valid. 

\begin{lemma}[Domain-inclusion]
\label{lemma-domain-inclusion}
Assume $R, S, T$ are any relations and $R_1, \ldots, R_n$ is any sequence of relations
on $[1,\MaxInd]$.
\begin{enumerate}

\item
\label{lemma-domain-inclusion-01}
For any two relations $R$, $S$ we have $\dom(RS) \subseteq \dom(R)$
and therefore $\DomCard(RS) \le \DomCard(R))$.

\item
\label{lemma-domain-inclusion-02}

$$
\dom(R_1) \supseteq 
\dom(R_1 R_2) \supseteq 
\dom(R_1 R_2 R_3) \supseteq 
\ldots \supseteq 
\dom(R_1 \ldots R_n)
$$

\item
\label{lemma-domain-inclusion-03}
$
(\dom(R) = \dom(RS))  
\ \ \  \Leftrightarrow \ \ \ 
(\dom(R) \subseteq \dom(RS)) 
\ \ \  \Leftrightarrow \ \ \  
(\dom(R) \not \supset  \dom(RS))
\ \ \  \Leftrightarrow \ \ \ 
(\DomCard(RS) = \DomCard(R)
$.

\item
\label{lemma-domain-inclusion-04}
If $\dom(S) = \dom(ST)$ then 
$\dom(RS) = \dom(RST)$.

\end{enumerate}
\end{lemma}

\begin{proof}
\begin{enumerate}

\item
%\label{lemma-domain-inclusion-01}
We have to prove that $\dom(RS) \subseteq \dom(R)$. For, we assume $x \in \dom(RS)$ in 
order to prove that $x \in \dom(R)$. If $x \in \dom(RS)$ then for some $z$ we have $x RS z$, 
therefore $x R y S z$ for some $y$. We conclude that $x \in \dom(R)$, as wished.

\item
%\label{lemma-domain-inclusion-02}
We prove 
$
\dom(R_1) \supseteq 
\dom(R_1 R_2) \supseteq 
\dom(R_1 R_2 R_3) \supseteq 
\ldots \supseteq 
\dom(R_1 \ldots R_n)
$ 
by induction on $n$. 
When $n=0$ there is no inclusion to prove, when $n=1$ we have 
$\dom(R_1) \supseteq \dom(R_1)$. Suppose $n \ge 1$ and the thesis 
for $n$, namely that $\dom(R_1) \supseteq \dom(R_1 R_2) \supseteq \dom(R_1 R_2 R_3) 
\supseteq \ldots \supseteq \dom(R_1 \ldots R_n)$. We have to prove the thesis for $n+1$. 
By point \ref{lemma-domain-inclusion-01}above 
we have $\dom(R_1 \ldots R_{n} R_{n+1}) \subseteq \dom(R_1 \ldots R_{n})$, which is
$\dom(R_1 \ldots R_{n}) \supseteq \dom(R_1 \ldots R_{n} R_{n+1})$. 
From this and the inductive hypothesis we deduce the thesis for $n+1$.

\item
%\label{lemma-domain-inclusion-03}
\begin{enumerate}
\item
\emph{Assume $(\dom(R) = \dom(RS))$}. 
We conclude $(\dom(R) \subseteq \dom(RS))$. 
\item
\emph{Assume $(\dom(R) \subseteq \dom(RS))$}. 
We conclude $(\dom(R) \not \supset \dom(RS))$.
\item
\emph{Assume $(\dom(R) \not \supset \dom(RS))$}. 
Then from $(\dom(R) \supseteq \dom(RS))$
(point \ref{lemma-domain-inclusion-01}above) we deduce that $(\dom(R) = \dom(RS))$.
We conclude $(\DomCard(R) = \DomCard(RS))$. 
\item
\emph{Assume $(\DomCard(R) = \DomCard(RS))$}, 
that is, $(\Card(\dom(R)) = \Card(\dom(RS)))$.
Both sets are included in $[1,\MaxInd]$ and therefore finite.
For finite sets, strict subsets have a smaller cardinality: if cardinality is the same
then an inclusion cannot be strict.
From $(\dom(R) \supseteq \dom(RS))$
(point \ref{lemma-domain-inclusion-01}above) 
and $(\Card(\dom(R)) = \Card(\dom(RS)))$ and finiteness 
we conclude that $(\dom(R) = \dom(RS))$.
\end{enumerate}

\item
%\label{lemma-domain-inclusion-04}
Assume that $\dom(S) = \dom(ST)$ in order to prove that 
$\dom(RS) = \dom(RST)$. 
We have $\dom(RS) \supseteq \dom(RST)$ by point \ref{lemma-domain-inclusion-01} above,
so we only have to prove that $\dom(RS) \subseteq \dom(RST)$.
For, we assume $a \in \dom(RS)$ in order to prove that $a \in \dom(RST)$. 
From $a \in \dom(RS)$ we deduce that there are $b \in \dom(S)$ and $c$ such that $a R b S c$. 
From  $\dom(S) = \dom(ST)$ we deduce that $b \in \dom(ST)$, namely that
for some $c' \in \dom(T)$ and $d$ we have $b S c' T d$. From $a R b S c' T d$ we conclude
that $a \in \dom(RST)$, as wished.

\end{enumerate}
\end{proof}

The inclusion $\dom(RS) \subseteq \dom(R)$ proved in
Lemma \ref{lemma-domain-inclusion}.\ref{lemma-domain-inclusion-01}
is often strict, as the next example says. 

\begin{example} 
\label{example-domain-inclusion}
$R = \{(1,1)\}$ and $S = \{(2,2)\}$ then $RS=\emptyset$,
therefore $\dom(RS) = \emptyset \subset \dom(R) = \{1\}$. 
\end{example}

We call domain-stable any sequence $\vec{S}$ of relations in which the strict
inclusion $\dom(S_i S_{i+1}) \subset \dom(S_i)$ never happens.
 
\begin{definition}[Domain-Stability]
Let $\vec{S}=S_1, \ldots, S_n$ be any sequence of relations. 

We say that $\vec{S}$ is domain-stable, only \emph{stable} for short, 
if and only if for all $i \in [1,n-1]$ we have $\dom(S_i S_{i+1}) = \dom(S_i)$.

\end{definition}


\begin{example}
\label{example-stable}
A sequence of length $n=0,1$ is trivially stable because $[1,n-1] = \emptyset$. 
Another example, with $n=3$. Let $R=\{(1,1)\}$, $S=\{(1,1),(2,2)\}$, $T=\{(1,1)\}$.
Then $\dom(ST) = \{1\} \subset \{1,2\} = \dom(S)$, therefore $R,S,T$ is \emph{not} stable.
Instead we have $R=T=RS=ST=RST=\{(1,1)\}$, therefore all other groupings
\begin{center}
$RS,T$ and $R,ST$ and $RST$
\end{center} 
are stable groupings.
\end{example}


By Lemma \ref{lemma-domain-inclusion}.\ref{lemma-domain-inclusion-03} we can
restate domain-stability in several ways. We have
$R,S$ stable $ \Leftrightarrow (\dom(R) = \dom(RS))  
  \Leftrightarrow  
(\dom(R) \subseteq \dom(RS)) 
  \Leftrightarrow   
(\dom(R) \not \supset  \dom(RS))
$.

Another way of saying that the pair $R,S$ is stable is that we do not remove elements
from $\dom(R)$ when we compose $R$ with $S$.

In particular, $R,S$ stable means that $(\dom(R) \subseteq \dom(RS))$, 
namely that whenever $x$ has an $R$-image $y$, then $x$ has some
$R$-image $y'$ having some $S$-image $z$. It does \emph{not} say, however,
that each $R$-image of $x$ has some $S$-image: it only says that if $x$ has some
$R$-image, then some of the $R$-images of $x$ have some $S$-image. 

In the same way we can restate the stability for a sequence.

\begin{proposition}[Restating Stability]
Let $\vec{S}=S_1, \ldots, S_n$ be any sequence of relations. 
Then  $\vec{S}$ is stable if and only if $S_i \cap (\dom(S_i) \times \dom(S_{i+1}))$ 
is everywhere defined for all $i \in [1,n-1]$. 
\end{proposition}

\begin{proof}
Indeed, this is just another way of saying that if $x$ has some $S_i$-image,
then some $S_i$-image of $x$ has some $S_{i+1}$-image. This is the definition
of $S_i$, $S_{i+1}$ domain-stable for $i \in [1,n-1]$. 
\end{proof}

We can restate stability in many more ways.





\begin{lemma}[Domain-Stability]
\label{lemma-domain-stability}
Let $\vec{S}=S_1, \ldots, S_n$ be any sequence of relations. The following are equivalent:

\begin{enumerate}
\item
\label{lemma-domain-stability-01}
$\vec{S}$ is stable

\item
\label{lemma-domain-stability-02}
for all $i, j \in [1, n]$, $i \le j$ we have:
$$
\dom(S_i) =  \dom(S_i \ldots S_j)
$$
or equivalently $\DomCard(S_i) =  \DomCard(S_i \ldots S_j)$.

\item
\label{lemma-domain-stability-03}
all grouping and all subsegments of $\vec{S}$ are stable.

\item 
\label{lemma-domain-stability-04}
for all $i \in [1,n-1]$ we have:
$$
\dom(S_i) = \dom(S_i \ldots S_n)
$$

\end{enumerate}
\end{lemma}



\begin{proof}
We can equivalently restate $\dom(S_i) =  \dom(S_i \ldots S_j)$ as
$\DomCard(S_i) = \DomCard(\dom(S_i \ldots S_j)$, 
given that we have $\dom(S_i) \supseteq  \dom(S_i \ldots S_j)$ 
by Lemma \ref{lemma-domain-inclusion}.\ref{lemma-domain-inclusion-02})
and the fact that they are subsets of the finite set $[1,\MaxInd]$, therefore 
strict subsets have a smaller cardinality.

So we can skip the formulation using 
$\DomCard(S_i) = \DomCard(\dom(S_i \ldots S_j)$
when deriving the equivalences about point \ref{lemma-domain-stability-02}.

We will now prove the equivalence of points 
\ref{lemma-domain-stability-01},
\ref{lemma-domain-stability-02},
\ref{lemma-domain-stability-03},
namely we prove that
\begin{center}
\ref{lemma-domain-stability-01} 
$\implies$ 
\ref{lemma-domain-stability-02}
$\implies$ 
\ref{lemma-domain-stability-03} 
$\implies$ 
\ref{lemma-domain-stability-01}
\end{center}
Then we prove the equivalence of points
\ref{lemma-domain-stability-02},
\ref{lemma-domain-stability-04},
by proving that 
\begin{center}
\ref{lemma-domain-stability-02} $\implies$  
\ref{lemma-domain-stability-04} $\implies$  
\ref{lemma-domain-stability-02}
\end{center}


\begin{enumerate}

\item
\ref{lemma-domain-stability-01} $\implies$ \ref{lemma-domain-stability-02}.

We assume point \ref{lemma-domain-stability-01}: $\vec{S}$ is stable, in to prove 
point \ref{lemma-domain-stability-02}: 
for all $i, j \in [1, n]$, $i \le j$ we have $\dom(S_i) =  \dom(S_i \ldots S_j)$.

For, we assume $j \ge i$ and we prove $\dom(S_i) =  \dom(S_i \ldots S_j)$ by induction on $j$.
\begin{enumerate}
\item
Assume $j=i$. Then $\dom(S_i) =  \dom(S_i)$.
\item
Assume $\dom(S_i) =  \dom(S_i \ldots S_j)$ and $j<n$ in order to prove that 
$\dom(S_i) =  \dom(S_i \ldots S_j S_{j+1})$. By definition of stability we have 
$\dom(S_j) = \dom(S_j S_{j+1})$. By Lemma 
\ref{lemma-domain-inclusion}.\ref{lemma-domain-inclusion-04} 
applied to $(S_i \ldots S_{j-1}), \ S_j, \ S_{j+1}$ we deduce 
$\dom( (S_i \ldots S_{j-1}) S_j) = \dom( (S_i \ldots S_{j-1}) S_j S_{j+1})$.
From the inductive hypothesis $\dom(S_i) =  \dom( (S_i \ldots S_{j-1}) S_j)$
we conclude that $\dom(S_i) =  \dom((S_i \ldots S_{j-1}) S_j S_{j+1})$, as wished.
\end{enumerate}


\item
\ref{lemma-domain-stability-02} $\implies$ \ref{lemma-domain-stability-03}.

We assume point \ref{lemma-domain-stability-02}: 
for all $i, j \in [1, n]$, $i \le j$ we have $\dom(S_i)  =  \dom(S_i \ldots S_j)$,
in order to prove point \ref{lemma-domain-stability-03}:
all subsegments and grouping of $\vec{S}$ are stable. 
\begin{enumerate}
\item
\emph{We prove that all grouping of $\vec{S}$ are stable}. We assume that
$i, j,k\in [1, n]$, $i \le j < k$ and that $(S_i \ldots S_j)$, $(S_{j+1} \ldots S_k)$
are two consecutive products in some grouping of $\vec{S}$, and we have to prove 
that $\dom(S_i \ldots S_j) = \dom( (S_i \ldots S_j) (S_{j+1} \ldots S_k))$. 
By point \ref{lemma-domain-stability-02} we have 
$$
\dom(S_i \ldots S_j) = \dom(S_i) = \dom(S_i \ldots S_j S_{j+1} \ldots S_k)
$$
\item
\emph{We prove that all sub-segments are stable}. 
We proved that all grouping of $\vec{S}$ are stable.
Therefore $\vec{S}$ itself is stable, being a trivial grouping of $\vec{S}$.
The finite quantification defining stability for a subsegment of $\vec{S}$ 
is a restriction of the finite quantification defining stability for $\vec{S}$. From
stability of $\vec{S}$ we deduce stability for each sub-segment.
\end{enumerate}
This proves point \ref{lemma-domain-stability-03}.


\item
\ref{lemma-domain-stability-03} $\implies$ \ref{lemma-domain-stability-01}.

We assume point \ref{lemma-domain-stability-03}: all sub-segments 
and groupings of $\vec{S}$ are stable,
and we have to prove point \ref{lemma-domain-stability-01}: $\vec{S}$ is stable. 

This follows from the fact that $\vec{S}$ itself is a trivial case of grouping of $\vec{S}$.


\item
\ref{lemma-domain-stability-02} $\implies$  \ref{lemma-domain-stability-04}.
We assume point \ref{lemma-domain-stability-02}:
for all $i, j \in [1, n]$, $i \le j$ we have $\dom(S_i) =  \dom(S_i \ldots S_j)$
and we have to prove point \ref{lemma-domain-stability-04}: 
for all $i \in [1, n-1]$ we have $\dom(S_i) =  \dom(S_i \ldots S_n)$.

This follows by choosing $j=n \ge i$.


\item
\ref{lemma-domain-stability-04} $\implies$  \ref{lemma-domain-stability-02}.

We assume point \ref{lemma-domain-stability-04}: 
for all $i \in [1, n-1]$ we have $\dom(S_i) =  \dom(S_i \ldots S_n)$, and we have to prove
point \ref{lemma-domain-stability-02}: 
for all $i, j \in [1, n]$, $i \le j$ we have $\dom(S_i) =  \dom(S_i \ldots S_j)$.

\emph{Suppose $i=n$}. 
Then the thesis becomes $\dom(S_n) =  \dom(S_n)$, which is trivially true.

\emph{Suppose $i \in [1,n-1]$}.
By Lemma \ref{lemma-domain-inclusion}.\ref{lemma-domain-inclusion-01},
for each sequence $\vec{S}$ of relations we have 
$$
\dom(S_i)  \ \ \ \supseteq \ \ \ 
\dom(S_iS_{i+1})  \ \ \  \supseteq \ \ \ 
\dom(S_iS_{i+1}S_{i+2})  \ \ \  \supseteq  \ \ \ 
\ldots  \ \ \  \supseteq  \ \ \ 
\dom(S_i \ldots S_n)
$$ 
therefore if we add the assumption $\dom(S_i \ldots S_n) = \dom(S_i)$ we conclude that
$$
\dom(S_i)   \ \ \ =  \ \ \ 
\dom(S_iS_{i+1})  \ \ \  =  \ \ \ 
\dom(S_iS_{i+1}S_{i+2})   \ \ \ =  \ \ \ 
\ldots   \ \ \ =  \ \ \ 
\dom(S_i \ldots S_n)
$$ 
as wished.
 
\end{enumerate}
\end{proof}




Lemma \ref{lemma-domain-stability}.\ref{lemma-domain-stability-02}
says that a sequence is stable if and only if whenever $x$ has some $S_i$-image,
then $x$ has some chain of $S_i, S_{i+1}, S_{i+2}, \ldots$-images of any length
moving through the sequence. We never drop any element of $\dom(S_i)$
when moving to the sets $\dom(S_i S_{i+1})$, $\dom(S_i S_{i+1} S_{i+2})$, \ldots.
This is an alternative definition of domain-stable sequence.





We claim that the grouping relation $\grouping$ is a partial order, but we include
no proof of this fact because it is not required for the paper. For instance, 
$S_1 \ldots S_n$ is the maximum grouping and $S_1, \ldots, S_n$ the minimum grouping.

All sequences $\vec{S}=S_1, \ldots, S_n$ have a trivial stable grouping, by taking the
product $S_1 \ldots S_n$ of the sequence. This is the maximum grouping w.r.t. $\grouping$, 
and therefore it is the maximum stable grouping, but it is trivial. 

There is no minimum
stable grouping w.r.t. $\grouping$. For instance,
in the Example \ref{example-stable} $R,ST$ and $RS,T$
are stable grouping and they are not comparable. The unique smaller grouping is 
$R,S,T$ and it is not stable.

However, we can extend the $\grouping$ relation to a total order $\le$:
by totality of $\le$ and finiteness of grouping of $\vec{S}$, there is a minimum stable grouping  w.r.t. $\le$ for each $\vec{S}$.
Namely, there is some stable $\vec{T} \grouping \vec{S}$ such that for all stable 
$\vec{U} \grouping \vec{S}$ we have $\vec{U} \le \vec{T}$. 

Lexicographic ordering compares two groupings, looking for the grouping
whose products which are closest to the left are shorter.

\begin{definition}[Lexicographic Ordering on Grouping]
\label{definition-grouping-order}
Assume $\vec{S}=S_1, \ldots, S_n$ is a sequence of relations and 
$$
\vec{R} \ \ \ 
= \ \ \ 
(S_{a_1} \ldots S_{b_1}), \ \ \ 
(S_{a_2} \ldots S_{b_2}), \ \ \ 
\ldots, \ \ \ 
(S_{a_{h}} \ldots S_{b_{h}})
$$
is a grouping of $\vec{S}$ with $a_1=1, a_{i+1} = b_i+1$ for all $i \in [1,h-1]$ and 
$b_{h} = n$.

\begin{enumerate}
\item
The degree of the grouping $\vec{R}$ of $\vec{S}$ is the sequence
of the lengths of its products, namely $(b_1, b_2 - b_1, \ldots, b_{h} - b_{h-1})$. 
\item
We compare two grouping of the same sequence by comparing their degrees in the 
lexicographic ordering.
\end{enumerate}
\end{definition} 

Lexicographic ordering $\le$ is a total order on a finite set, and there is the trivial stable
grouping consisting of a single product,
therefore there is a minimum stable grouping w.r.t. $\le$.


We provide a characterization of the minimum stable grouping of 
$\vec{S}$ using a map $\phi(\vec{S},\cdot) : [1,\MaxInd] \rightarrow  [1,\MaxInd]$. 
Then using the map $\phi(\vec{S},\cdot)$ we will provide an algorithm computing 
the minimum stable group in polynomial time. 
We first need the notion of $j$ dropping index for $i$, for $i,j \in [1,n]$.

\begin{definition}[Dropping Index]
\label{definition-dropping-index}
Let $\vec{S} = S_1, \ldots, S_n$ and $i,j \in [1,n]$, $i \le j$.
We say that $j$ is a dropping index for $i$ in $\vec{S}$
if and only if $j=i$ or $\dom(S_i \ldots S_{j-1}) \supset \dom(S_i \ldots S_j)$.  
\end{definition}

By definition unfolding, $j$ dropping index for $i$ in $\vec{S}$
if and only if either $j=i$ or $j>i$ and there is some $x \in [1,\MaxInd]$ having some 
$S_i \ldots S_{j-1}$-image  but no $S_i \ldots S_j$-image: 
some traces from $x$ continue until step $j-1$, then they all die out. 


\begin{definition}[Maximum Dropping index $\phi(\vec{S},i)$]
Let $\vec{S} = S_1, \ldots, S_n$ and $i \in [1,n]$.
$\phi(\vec{S},i) \in [i,n]$ is the maximum dropping index $j\in [i,n]$ 
of $i$ in $\vec{S}$. 

We write just $\phi(i)$ when $\vec{S}$ is clear from the context. 
\end{definition}

If $\vec{S}$ has $n$ elements, then $\phi(\vec{S},n) \in [n,n]$ and therefore
$\phi(\vec{S},n)=n$.

In the Example \ref{example-stable}, with  $n=3$, $R=\{(1,1)\}$, $S=\{(1,1),(2,2)\}$, 
$T=\{(1,1)\}$, we have $\phi(1)=1$ because $\dom(R) = \dom(RS) = \dom(RST) =  \{1\}$.
We have $\phi(2)=3$ because $\dom(ST) = \{1\} \subset \{1,2\} = \dom(S)$. 
From $\phi(n)=n$ and $n=3$ we deduce $\phi(3)=3$.

\begin{lemma}[Equivalent Definition of $\phi$]
\label{lemma-equivalent-phi}
Assume $\vec{S} = S_1, \ldots, S_n$ is any relation sequence and $i,j \in [i,n]$.
Let us abbreviate $\phi(i) = \phi(\vec{S},i)$.

\begin{enumerate}
\item
\label{lemma-equivalent-phi-01}
$\dom(S_i \ldots S_{\phi(i)}) = \dom(S_i \ldots S_n)$.

\item
\label{lemma-equivalent-phi-02}
$\phi(i)$ is the minimum $j \in [i,n]$ such that
$\dom(S_i \ldots S_j) = \dom(S_i \ldots S_n)$.

\item
\label{lemma-equivalent-phi-03}
If $S_a \ldots S_b$ is a product in a stable grouping of $\vec{S}$, then 
$\phi(\vec{S},a) \le b$.

\end{enumerate}

\end{lemma}

\begin{proof}

\begin{enumerate}

\item
$\phi(i)$ is the last dropping index for $i$. By definition of dropping index we have
$\dom(S_i \ldots S_{k-1}) = \dom(S_i \ldots S_{k})$ for all $k \in [\phi(i)+1,n]$. By induction
on $k \in [\phi(i),n]$ we deduce that 
$\dom(S_i \ldots S_{\phi(i)}) = \dom(S_i \ldots S_k)$. We conclude that 
$\dom(S_i \ldots S_{\phi(i)}) = \dom(S_i \ldots S_n)$.

\item
Assume $j \in [i,n]$ is the first index such that $\dom(S_i \ldots S_{j}) = \dom(S_i \ldots S_n)$. 
By the previous point we have $j \le \phi(i)$.

By Lemma \ref{lemma-domain-inclusion}.\ref{lemma-domain-inclusion-02}
for all $k \in [j+1,n]$ we have 
$\dom(S_i \ldots S_j) \supseteq \dom(S_i \ldots S_{k-1}) \supseteq 
\dom(S_i \ldots S_{k}) \supseteq \dom(S_i \ldots S_n)$.
If we add $\dom(S_i \ldots S_{j}) = \dom(S_i \ldots S_n)$ we deduce 
$\dom(S_i \ldots S_j) = \dom(S_i \ldots S_{k-1}) = 
\dom(S_i \ldots S_{k}) = \dom(S_i \ldots S_n)$, namely that there is no dropping
index $k$ of $i$ with $k>j$. This means that $j \ge \phi(i)$. 

We conclude that $j = \phi(i)$.

\item
\label{lemma-equivalent-phi-03}
Assume that $\vec{T} = T_1, \ldots, T_m$ is a stable grouping of $\vec{S}$
with $T_i= S_{a_i} \ldots S_{b_i}$, in order to prove that $\phi(a_i) \le b_i$. 
By Lemma \ref{lemma-domain-stability}.\ref{lemma-domain-stability-02} 
and stability of $\vec{T}$ we have $\dom(T_i) = \dom(T_i \ldots T_m)$,
that is, $\dom(S_{a_i} \ldots S_{b_i}) = \dom(S_{a_i} \ldots S_{b_i} S_{b_i +1} \ldots S_n)$.
By point \ref{lemma-equivalent-phi-02} above, this implies $\phi(\vec{S},a_i) \le b_i$.

\end{enumerate}
\end{proof}

%By definition unfolding, $i \in [1,n]$ is $\phi$-closed if and only if for all $a \in [1,i[$,
%for all dropping index $b \in [1,n]$ of $a$ we have $b \in [1,i[$.

We define a grouping $\vec{S}^{\phi}$ of $\vec{S}$ using the map $\phi$,
then we prove that it stable and the minimum stable.

\begin{definition} [A Grouping Defined by the Map $\phi$]
Assume $\vec{S}$ is any relation sequence.
Let $a_1 < \ldots < a_h$ be the unique integer sequence such that:
$a_1=1$ and $a_{i+1} = \phi(a_i)+1$ for all $i \in [1,h-1]$ and $\phi(a_{h}) = n$.
We define $\vec{S}^{\phi} = R_1, \ldots, R_h$ by: 
$$
R_i = (S_{a_i}\ldots S_{\phi(a_i)}) \ \ \ \mbox{ for all $i \in [1,h]$}
$$ 
\end{definition}

We first check that $\vec{S}^{\phi}$ is a stable grouping of $\vec{S}$, 
then that $\vec{S}^{\phi}$ is the minimum stable.

\begin{lemma} [A Stable Grouping Defined by the Map $\phi$]
\label{lemma-stable-phi-1}
$\vec{S}^\phi$ is a stable grouping of $\vec{S}$.
\end{lemma}

\begin{proof}
Let $\vec{S}^\phi = R_1, \ldots, R_h$. For all $i \in [1,h]$ we have
$
\dom(R_i) = 
$ 
(by definition of $R_i$)  = 
$
\dom(S_{a_i}\ldots S_{\phi(a_i)}) = 
$ 
(by Lemma \ref{lemma-equivalent-phi})  = 
$
\dom(S_{a_i}\ldots S_{n}) =
$ 
(by definition of $R_i, \ldots, R_h$)  = 
$
\dom(R_{i}\ldots R_h)
$.
We deduce that the grouping $\vec{R}$ is stable by 
Lemma \ref{lemma-domain-stability}.\ref{lemma-domain-stability-04}.
\end{proof}

\begin{lemma}[Minimum Stable Grouping]
\label{lemma-stable-phi-2}
$\vec{S}^\phi$ is the minimum stable grouping of $\vec{S}$
w.r.t. the lexicographic ordering.
\end{lemma}

\begin{proof}
Assume $S_{a_1}, \ldots, S_{b_1}$ is the first product of a stable
grouping (hence $a_1=1$). 
By Lemma \ref{lemma-equivalent-phi}.\ref{lemma-equivalent-phi-03}.
we have $b_1 \ge \phi(a_1)$, hence the shortest first product is 
$S_{a_1} \ldots S_{b_1}$, with $a_1 = 1$ and $b_1 = \phi(a_1)$.
By the same reasoning, 
for any stable grouping starting with the product $S_{a_1} \ldots S_{b_1}$, 
the shortest second group is $S_{a_2}, \ldots, S_{b_2}$, 
with $a_2 = b_1 + 1$ and $b_2 = \phi(a_2)$. By induction on the number of elements 
of $\vec{S}^\phi$ we can prove that if $\vec{T} \grouping \vec{S}$ is stable then
$\vec{T}  \ge \vec{S}^\phi$ lexicographically.
\end{proof}


\begin{definition}[The Minimum Stable Grouping $\minStable(\vec{S})$]
\label{definition-minstable}
Assume $\vec{S}$ is any sequence of relations. 
\begin{enumerate}
\item
$\minStable(\vec{S}) = \vec{S}^\phi$
is the minimum stable grouping of $\vec{S}$ w.r.t. the lexicographic ordering $\le$.
\item
$(S_{a_h} \ldots S_{b_h})$, the last product of $\minStable(\vec{S}) $, 
is the stable tail of $\vec{S}$.
\end{enumerate}
\end{definition}


We are looking for a polynomial time algorithm for $\minStable(\vec{S})$.
$\phi(\vec{S},i)$ can be computed in polynomial time
by induction on $\vec{S}$, as follows. 

\begin{lemma}[Inductive Characterization of $\phi(\vec{S},i)$]
\label{lemma-phi-characterization}
We have $\phi(\vec{S} S_{n+1}, n+1) = n+1$ and for all $i \in [1,n]$
one of the following holds:
\begin{enumerate}
\item
either $\dom(S_i \ldots S_n) = \dom(S_i \ldots S_n S_{n+1})$
and $\phi(\vec{S}S_{n+1},i) = \phi(\vec{S},i)$;
\item
or $\dom(S_i \ldots S_n) \supset \dom(S_i \ldots S_n S_{n+1})$
and $\phi(\vec{S}S_{n+1},i) = n+1$.
\end{enumerate}
$\phi(\vec{S},\cdot):[1,n] \rightarrow [1,n]$ is a weakly increasing map.
\end{lemma}

\begin{proof}
$\phi(\vec{S} S_{n+1}, n+1) = n+1$ follows from $\phi(\vec{S} S_{n+1}, n+1) \ge n+1$
and $\phi(\vec{S} S_{n+1}, n+1) \in [1,n+1]$. 
By Lemma \ref{lemma-domain-inclusion}.\ref{lemma-domain-inclusion-01}, 
either $\dom(S_i \ldots S_n) = \dom(S_i \ldots S_n S_{n+1})$ 
or $\dom(S_i \ldots S_n) \supset \dom(S_i \ldots S_n S_{n+1})$. 
We argue by cases.

\begin{enumerate}
\item
\emph{Suppose $\dom(S_i \ldots S_n) = \dom(S_i \ldots S_n S_{n+1})$}.
By Lemma \ref{lemma-equivalent-phi}.\ref{lemma-equivalent-phi-02} 
for $j = \phi(\vec{S},i)$
we have $\dom(S_i \ldots S_j) = \dom(S_i \ldots S_n) = \dom(S_i \ldots S_n S_{n+1})$.
By Lemma \ref{lemma-equivalent-phi}.\ref{lemma-equivalent-phi-02} 
for all $k<j$ we have 
$\dom(S_i \ldots S_k) \not = \dom(S_i \ldots S_n) = \dom(S_i \ldots S_n S_{n+1})$.
Then Lemma \ref{lemma-equivalent-phi}.\ref{lemma-equivalent-phi-02} 
implies $\phi(\vec{S}S_{n+1},i) = \phi(\vec{S},i)$

\item
\emph{Suppose $\dom(S_i \ldots S_n) \supset \dom(S_i \ldots S_n S_{n+1})$}.
Then $n+1$ is a dropping index for $i$ in $\vec{S} S_{n+1}$.
We deduce $\phi(\vec{S}S_{n+1}, i) \ge n+1$,
namely $\phi(\vec{S}S_{n+1}, i) = n+1$.

\end{enumerate}

Now we check that $\phi$ is weakly increasing.

We assume $i,j \in [1,n]$, $i < j$ in order to prove that $\phi(\vec{S},i) \le \phi(\vec{S},j)$.
Let $b =  \phi(\vec{S},j)$. 
Then by Lemma \ref{lemma-equivalent-phi}.\ref{lemma-equivalent-phi-02} we have 
$\dom(S_j \ldots S_b) = \dom(S_j \ldots S_n)$. 
By Lemma \ref{lemma-domain-inclusion}.\ref{lemma-domain-inclusion-04} we deduce
$\dom((S_i \ldots S_{j-1}) S_j \ldots S_b) = \dom((S_i \ldots S_{j-1}) S_j \ldots S_n)$.
By definition of $\phi$, this implies $\phi(\vec{S},i) \le b =  \phi(\vec{S},j)$,
as wished.
\end{proof}

Using the characterization of $\phi(\vec{S},i)$, we have a polynomial
way of computing $\minStable(\vec{S})$ by induction on $\vec{S}$.

\begin{lemma}[Computing $\minStable(\vec{S})$ by induction on $\vec{S}$]
\label{lemma-minstable-inductive}
Assume $\vec{S} = S_1, \ldots, S_n$. 
\begin{enumerate}
\item
If $n=0$, then $\minStable(\vec{S}) = \vec{S} = $ the empty sequence.
\item
Suppose $\minStable(\vec{S}) = \vec{T} = T_1, \ldots, T_m$. Then exactly
one of the following holds.
\begin{enumerate}
\item
$\minStable(\vec{S}), S_{n+1}$ is stable and 
$\minStable(\vec{S}, S_{n+1}) = \minStable(\vec{S}), S_{n+1}$.
\item
$\minStable(\vec{S}), S_{n+1}$ is \emph{not} stable and
$\minStable(\vec{S}, S_{n+1}) = T_1, \ldots, (T_i \ldots T_m S_{n+1})$, 
for the first $i \in [1,m]$ such that 
$\dom(T_i \ldots T_m) \supset \dom(T_i \ldots T_m S_{n+1})$.
\end{enumerate}
\end{enumerate}
\end{lemma}

\begin{proof}
\begin{enumerate}
\item
If $n=0$, then $\minStable(\vec{S}) = \vec{S} = $ the empty sequence by definition
of $\minStable$.

\item
Suppose $\minStable(\vec{S}) = \vec{T} = T_1, \ldots, T_m$.  
By Lemma \ref{lemma-stable-phi-2}
we have $T_i = S_{a_i} \ldots S_{\phi(\vec{S},a_i)}$ with 
$a_1=1$ and $a_{i+1} = \phi(\vec{S},a_i)+1$ 
for all $i \in [1,m-1]$ and $\phi(a_{m}) = n$.  Then exactly one of the following holds.

\begin{enumerate}

\item
Assume $\minStable(\vec{S}), S_{n+1}$ is stable. By 
Lemma \ref{lemma-domain-stability-02}.
\ref{lemma-domain-stability-02},\ref{lemma-domain-stability-04} 
we have
$\dom(T_i \ldots T_m) = \dom(T_i ) = \dom(T_i \ldots T_m S_{n+1})$ for all $i \in [1,m]$,
namely $\dom(S_{a_i} \ldots S_n) = \dom(S_{a_i} \ldots  S_{n+1})$ for all $i \in [1,m]$.
From Lemma \ref{lemma-phi-characterization} we deduce that $\phi(\vec{S},a_i) = 
\phi(\vec{S}S_{n+1},a_i)$ for all $i \in [1,m]$. 

By induction on $i \in [1,m]$ we deduce that the first $m$ of products of 
$\minStable(\vec{S},S_{n+1})$ are $S_{a_i} \ldots S_{\phi(\vec{S},S,a_i)}$, which is
$S_{a_i} \ldots S_{\phi(\vec{S},a_i)}$, namely $T_i$.

We also have  $\phi(\vec{S},S_{n+1},\phi(a_{m}) +1)
= \phi(\vec{S},S_{n+1},n+1)=n+1$. We deduce that the last product is 
the unary product $S_{n+1}$.

We conclude that $\minStable(\vec{S}, S_{n+1}) = \minStable(\vec{S}), S_{n+1}$.

\item
Assume $\minStable(\vec{S}), S_{n+1}$ is \emph{not} stable. By 
Lemma \ref{lemma-domain-stability-02}.
\ref{lemma-domain-stability-02},\ref{lemma-domain-stability-04} 
we have
$\dom(T_i \ldots T_m) = \dom(T_i ) \not = \dom(T_i \ldots T_m S_{n+1})$ 
for a first $i \in [1,m]$.
By Lemma \ref{lemma-domain-inclusion}.\ref{lemma-domain-inclusion-01}
we have $\dom(T_i ) \supseteq  \dom(T_i \ldots T_m S_{n+1})$:
we deduce $\dom(T_i \ldots T_m) \supset \dom(T_i \ldots T_m S_{n+1})$.

By the choice of $i$ that 
$\dom(T_j \ldots T_m) = \dom(T_j ) = \dom(T_j \ldots T_m S_{n+1})$ for all $j \in [1,i-1]$.

By Lemma \ref{lemma-phi-characterization} we have 
$\phi(\vec{S},S_{n+1},a_j) = \phi(\vec{S},a_j)$ for all $j \in [1,i-1]$, while 
$\phi(\vec{S},S_{n+1},a_i) = n+1$. 

By Lemma \ref{lemma-stable-phi-2} and induction on $j \in [1,i-1]$,
the product number $j$ of $\minStable(\vec{S},S_{n+1})$ is $T_j$.
Indeed, assume that $j=1$ or $j>1$ and the previous product is $T_{j-1}$.
In both cases, by Lemma \ref{lemma-stable-phi-2} the first element
of the product number $j$ of $\minStable(\vec{S},S_{n+1})$
is $a_j$, with $a_j = 1$ if $j=1$, and $a_j = \phi(\vec{S},a_ {j-1}+1)$ if $j>1$. 
Again by by Lemma \ref{lemma-stable-phi-2} the product number $j$ 
of $\minStable(\vec{S},S_{n+1})$ is $S_{a_j} \ldots S_{\phi(\vec{S},S_{n+1},a_j)}$, 
which is $S_{a_i} \ldots S_{\phi(\vec{S},a_j)}$, namely $T_j$.
The product number $i$ is the last product and it is  
$S_{a_i} \ldots S_{\phi(\vec{S},S_{n+1},a_i)}$,
namely $S_{a_i} \ldots S_{n+1}$.

We conclude that 
$\minStable(\vec{S}, S_{n+1}) = T_1, \ldots, T_{i-1}, (T_i \ldots T_m S_{n+1})$.
\end{enumerate}
\end{enumerate}
\end{proof}


\section{Injective Domain-Stable Relation Sequences}

We first derive some inequalities for domain-cardinality $\DomCard(\vec{S})$
valid for sequences of \emph{injective} relations.

\begin{lemma}[Injective Relation Sequences]
\label{lemma-injective-sequences}
Assume $R$, $S$ and $\vec{R}=R_1, \ldots, R_n$ all are \emph{injective} relations.

\begin{enumerate}
\item
\label{lemma-injective-sequences-01}
$\DomCard(RS) \le \DomCard(S)$
and if $\DomCard(RS) = \DomCard(S)$ then $R$ restricted to $\dom(R) \times \dom(S)$
is a bijection.

\item
\label{lemma-injective-sequences-02}
For all $i,j \in [1,n]$, $i \le j$:
$$
\DomCard(R_i \ldots R_n) \le \DomCard(R_j \ldots R_n)
$$

\item
\label{lemma-injective-sequences-03}
Assume $i, j \in[1,n]$ and $i < j$ and 
$\Card(\dom(R_i \ldots R_n)) = \Card(\dom(R_j \ldots R_n))$.
Then $R_i \ldots R_{j-1}$ restricted to $\dom(R_i \ldots R_n) \times \dom(R_j \ldots R_n)$
is a bijection.
\end{enumerate}
\end{lemma}

\begin{proof}

Assume $R, S$, $\vec{R}=R_1, \ldots, R_n$ all are \emph{injective} relations.

\begin{enumerate}

\item
%\label{lemma-injective-sequences-01}
For all $x \in \dom(RS)$ the sets $\dom(S) \cap R(x)$ are inhabited,
because if $x \in \dom(RS)$ then for some $y \in \dom(S), z$ we have $x R y S z$.
They are pairwise disjoint subsets of $\dom(S)$, because $R$ is injective. We deduce that 
$$
\DomCard(RS) = 
\Card(\dom(RS)) \le
\Card(\bigcup_{x \in \dom(RS)} \dom(S) \cap R(x)) \le 
\Card(\dom(S)) = 
\DomCard(S)
$$

Now assume that $\DomCard(RS) = \DomCard(S)$. Then 
$$
\DomCard(RS) = \Card(\bigcup_{x \in \dom(RS)} \dom(S) \cap R(x)) = \DomCard(S)
$$ 
From 
$$
\Card(\bigcup_{x \in \dom(RS)} \dom(S) \cap R(x)) = \Card(\dom(S)), \ \ \ 
\bigcup_{x \in \dom(RS)} \dom(S) \cap R(x) \subseteq \dom(S)
$$ 
and $\dom(S)$ finite 
we deduce that $$\bigcup_{x \in \dom(RS)} \dom(S) \cap R(x) = \dom(S)$$ 
From the fact that $\Card(\dom(RS)) = \Card(\dom(S))$ and all $\dom(S) \cap R(x)$ are 
inhabited we deduce that all $\dom(S) \cap R(x)$ are singleton. 
This says that $R$ restricted to $\dom(R) \times \dom(S)$ is a bijection.

\item
%\label{lemma-injective-sequences-02}
Let $i \in [1,n-1]$.
From the point 1 above applied to $R = R_i$, $S = R_{i+1}\ldots R_n$ we deduce that
$\Card(\dom(R_i R_{i+1}\ldots R_n)) \le \Card(\dom(R_{i+1} \ldots R_n))$. The thesis
follows by induction on $j-i$.

\item
%\label{lemma-injective-sequences-03}
Assume $i, j \in[1,n]$ and $i < j$ and 
$\Card(\dom(R_i \ldots R_n)) = \Card(\dom(R_j \ldots R_n))$.
From point \ref{lemma-injective-sequences-01} 
above applied to $R = R_i \ldots R_{j-1}$ and $S = R_i \ldots R_{j-1}$ 
we deduce that $R_i \ldots R_{j-1}$ restricted to 
$\dom(R_i \ldots R_{j-1} \ldots R_j \ldots R_n) \times \dom(R_i \ldots R_{j-1})$
is a bijection.

\end{enumerate}
\end{proof}

We derive more inequalities for domain-cardinality
only valid for \emph{stable} sequences of injective relations. 


\begin{lemma}[Injective Domain-Stable Relation Sequences]
\label{lemma-injective-stable-sequences}
Assume $\vec{R}=R_1, \ldots, R_n$ is a \emph{domain-stable}
sequence of \emph{injective} relations.
\begin{enumerate}

\item
\label{lemma-injective-stable-sequences-01}
$$
\forall i,j \in [1,n]. (i \le j) \implies \DomCard(R_i) \le \DomCard(R_j)
$$

\item
\label{lemma-injective-stable-sequences-02}
If  $i, j \in[1,n]$ and $i < j$ and $\DomCard(R_i) = \DomCard(R_j)$
then $R_i \ldots R_{j-1}$ has domain $\dom(R_i)$ and
when restricted to $\dom(R_i) \times \dom(R_j)$
is a bijection.
\end{enumerate}
\end{lemma}


\begin{proof}
\begin{enumerate}

\item
%\label{lemma-injective-stable-sequences-01}
Assume $i,j \in [1,n]$, $i \le j$. By stability and
Lemma \ref{lemma-domain-stability}.\ref{lemma-domain-stability-04}
we have $\dom(R_i) = \dom(R_i \ldots R_n)$ and  $\dom(R_j) = \dom(R_j \ldots R_n)$:

We deduce $\DomCard(R_i) = \DomCard(R_i \ldots R_n))$ and
$\DomCard(R_j) = \DomCard(R_j \ldots R_n))$. By injectivity and 
Lemma \ref{lemma-injective-sequences}.\ref{lemma-injective-sequences-02}  
we obtain $\DomCard(R_i \ldots R_n) \le \DomCard(R_j \ldots R_n)$.

We conclude that $\DomCard(R_i) \le \DomCard(R_j)$.

\item
%\label{lemma-injective-stable-sequences-02}
Assume $i,j \in [1,n]$, $i < j$ and $\DomCard(R_i) = \DomCard(R_j)$.

By stability and
Lemma \ref{lemma-domain-stability}.\ref{lemma-domain-stability-04}
we have $\dom(R_i) = \dom(R_i \ldots R_{j-1})$, this is the first half of our thesis.

From the same Lemma \ref{lemma-domain-stability}.\ref{lemma-domain-stability-04}
we deduce $\dom(R_i) = \dom(R_i \ldots R_n)$ and $\dom(R_j) = \dom(R_j \ldots R_n)$. 
Thus, $\DomCard(R_i) = \DomCard(R_i \ldots R_n)$ and
$\DomCard(R_j) = \DomCard(R_j \ldots R_n)$. 
By assumption we have $\DomCard(R_i) = \DomCard(R_j)$,
we deduce $ \DomCard(R_i \ldots R_n) = \DomCard(R_j \ldots R_n)$.

By injectivity and Lemma \ref{lemma-injective-sequences-03} we obtain: 
$R_i \ldots R_{j-1}$ restricted to $\dom(R_i \ldots R_n) \times \dom(R_j \ldots R_n)$
is a bijection. 

By $\dom(R_i) = \dom(R_i \ldots R_n)$ and 
$\dom(R_j) = \dom(R_j \ldots R_n)$ we conclude that
$R_i \ldots R_{j-1}$ restricted to $\dom(R_i) \times \dom(R_j)$
is a bijection. 
\end{enumerate}

\end{proof}

%11/02/2026 11:24
We define a property stronger than domain-stability: having constant domain-cardinality
for all product of non-empty sub-segments of $\vec{S}$.


\begin{definition}[Constant Domain-Cardinality]
Assume $\vec{S}=S_1, \ldots S_n$ is any sequence of relations.
\begin{enumerate}

\item
$\vec{S}$ is a constant domain-cardinality sequence if and only if $\vec{S}$ is is stable
and for all $i,j \in [1,n]$:
$$
\DomCard(S_i) = \DomCard(S_j)
$$

\item
If $\vec{S}$ is not empty, then we call $\DomCard(\vec{S}) = \DomCard(S_1 \ldots S_n)$ 
the domain-cardinality of $\vec{S}$.

\item
A \emph{component} $\vec{C}$ of $\vec{S}$ is a maximal subsegment 
in $\vec{S}$ which is a constant domain-cardinality sequence.
\end{enumerate}
\end{definition}

If $\vec{S}=S_1, \ldots S_n$ is constant domain-cardinality, then by definition
$\DomCard(S_i) = \DomCard(S_j)$ for all $i,j \in [1,n]$. 

For all $k \in [i,n]$, $h \in [j,n]$ 
by Lemma \ref{lemma-domain-stability-02}.\ref{lemma-domain-stability-02} 
we have $\DomCard(S_i) =\DomCard(S_i \ldots S_k)$ and
$\DomCard(S_j) = \Card(\dom(S_j \ldots S_h))$. This means that all products of
all non-empty segments $[i,j],[k,h] \subseteq [1,n]$ have domains of the same cardinality: 
$\DomCard(S_i \ldots S_k) = \DomCard(S_j \ldots S_h)$.

By Lemma \ref{lemma-injective-sequences}.\ref{lemma-injective-sequences-01}, 
for all $i \in [1,n-1]$, $S_i$ is a bijection from $\dom(S_i)$ to
$\dom(S_{i+1})$.  We can picture the traces of $\vec{S}$ as a composition of 
bijections from $\dom(S_i)$ to $\dom(S_{i+1})$, 
plus some traces passing through $\dom(S_i)$, missing $\dom(S_{i+1})$ and dying out.

\begin{center}
\begin{tikzpicture}
\draw[->,black, thick] (0,0) -- (1,-1);
\draw[->,black, thick] (0,0) -- (1,0);
\draw[->,black, thick] (0,1) -- (1,1);
\draw[->,black, thick] (0,2) -- (1,2);

\draw[->,black, thick] (1,0) -- (2,-1);
\draw[->,black, thick] (1,0) -- (2,0);
\draw[->,black, thick] (1,1) -- (2,1);
\draw[->,black, thick] (1,2) -- (2,2);

\draw[->,black, thick] (2,0) -- (3,-1);
\draw[->,black, thick] (2,0) -- (3,0);
\draw[->,black, thick] (2,1) -- (3,1);
\draw[->,black, thick] (2,2) -- (3,2);

\filldraw[black] (0,0) circle (1pt) node[anchor=north]{$a_1$};
\filldraw[black] (0,1) circle (1pt) node[anchor=north]{$b_1$};
\filldraw[black] (0,2) circle (1pt) node[anchor=north]{$c_1$};

\filldraw[black] (1,-1) circle (1pt) node[anchor=north]{};
\filldraw[black] (1,0) circle (1pt) node[anchor=north]{$a_2$};
\filldraw[black] (1,1) circle (1pt) node[anchor=north]{$b_2$};
\filldraw[black] (1,2) circle (1pt) node[anchor=north]{$c_2$};

\filldraw[black] (2,-1) circle (1pt) node[anchor=north]{};
\filldraw[black] (2,0) circle (1pt) node[anchor=north]{$a_3$};
\filldraw[black] (2,1) circle (1pt) node[anchor=north]{$b_3$};
\filldraw[black] (2,2) circle (1pt) node[anchor=north]{$c_3$};

\filldraw[black] (3,-1) circle (1pt) node[anchor=north]{};
\filldraw[black] (3,0) circle (1pt) node[anchor=north]{$\ldots$};
\filldraw[black] (3,1) circle (1pt) node[anchor=north]{$\ldots$};
\filldraw[black] (3,2) circle (1pt) node[anchor=north]{$\ldots$};

\node[] at (1,3.2) {\tiny $\cod(S_1)$};
\node[] at (2,3.2) {\tiny $\cod(S_2)$};
\node[] at (3,3.2) {\tiny $\cod(S_3)$};

\node[] at (0,2.6) {\tiny $\dom(S_1)$};
\node[] at (1,2.6) {\tiny $\dom(S_2)$};
\node[] at (2,2.6) {\tiny $\dom(S_3)$};
\node[] at (3,2.6) {\tiny $\ldots$};

\draw[blue] (0,1) ellipse (.5cm and 1.5cm);
\draw[blue] (1,1) ellipse (.5cm and 1.5cm);
\draw[blue] (2,1) ellipse (.5cm and 1.5cm);
\draw[blue] (3,1) ellipse (.5cm and 1.5cm);
  
\draw[red] (1,1) ellipse (.5cm and 2.1cm);
\draw[red] (2,1) ellipse (.5cm and 2.1cm);
\draw[red] (3,1) ellipse (.5cm and 2.1cm);
\end{tikzpicture}
\end{center}

Each sequence can be transformed into a stable sequence by a minimum of grouping,
then into a concatenation of constant domain-cardinality sequences with increasing
domain-cardinality.

%18:22 04/02/2026

\begin{lemma}[Decomposing Stable Sequences of Injective Relations]
\label{lemma-decomposition-stable}
Every stable sequence $\vec{S}$ of injective relations is equal to the 
concatenation $\vec{C_1}, \ldots \vec{C_h}$
of some non-empty constant domain-cardinality subsegments 
such that $\DomCard(\vec{C_i})  < \DomCard(\vec{C_j})$ for all $i,j \in [1,h]$, $i < j$.
\end{lemma}


\begin{proof}
We define $a_1=1$, $b_i = $ the last $j \in [i,n]$ such that 
$\DomCard(S_{a_i}) = \DomCard(S_j)$, and if $b_i<n$ then $a_{i+1}=b_i + 1$,
otherwise the sequence $a_1, \ldots, a_i$ stops.
Then for some $h \in [1,n]$ we have:

$$
\vec{S}
=
S_{a_1},\ldots,S_{b_1},
\ldots
S_{a_{h}},\ldots,S_{b_{h}}
$$

Each $S_{a_{i}},\ldots,S_{b_{i}}$ is a sub-segment of $\vec{S}$, therefore it is stable
by Lemma \ref{lemma-domain-stability}.\ref{lemma-domain-stability-04}, 
and all sets $\dom(S_{a_{i}}), \ldots, \dom(S_{b_{i}})$ have the same cardinality 
by the choice of $b_i$. 
Thus, each $S_{a_{i}},\ldots,S_{b_{i}}$ has constant domain-cardinality. 

By Lemma \ref{lemma-injective-stable-sequences}.\ref{lemma-injective-stable-sequences-01}
we have $\DomCard(S_{a_i}) \le \DomCard(S_{a_{i+1}})$ for all $i \in [1,h-1]$.
By the choice of $b_i$ we also have 
$\DomCard(S_{a_{i}}) = \DomCard(S_{b_{i}}) \not = \DomCard(S_{a_{i+1}})$.
We deduce that
$\DomCard(S_{a_i}) <  \DomCard(S_{a_{i+1}})$ for all $i \in [1,h-1]$.
We conclude  that
$\DomCard(S_{a_i}) <  \DomCard(S_{a_{j}})$ for all $i,j \in [1,h]$, $i < j$,
by induction on $j$.

\end{proof}

Remark that the subsequence $S_{a_{i}}, \ldots, S_{b_{i}}$ 
are exactly the \emph{maximal} constant domain-cardinality subsegments of $\vec{S}$,
what we just called the \emph{components} of $\vec{S}$.
We call the set of indexes $a_i$ the \emph{raising indexes} of $\vec{S}$,
and the suffixes $S_{a_i} \ldots S_n$ raising suffixes (those in which the domain-cardinality
raises). The notion of raising indexes can be extended to any sequence. 


\begin{definition}[Raising Suffixes of $\vec{S}$] 
Assume $\vec{S}$ is sequence of injective relations. 
\begin{enumerate}

\item
Assume $\vec{S}$ is \emph{stable}. We call \emph{raising index} the first index $a_i$ of 
$S_{a_{i}}, \ldots, S_{b_{i}}$,
the $i$-th component of $\vec{S}$, with $i \in [1,h]$.
We define the set of \emph{raising suffixes} 
$$
\Raising(\vec{S}) = \{S_{a_i} \ldots S_n | i \in [1,h] \}
$$

\item
In general we set 
$$
\Raising(\vec{S}) = \Raising(\minStable(\vec{S}))
$$
\end{enumerate}
\end{definition}

%By definition unfolding, $\Raising(\vec{S})$ consists of all products $S_{a_i} \ldots S_n$ 
%of suffixes of $\vec{S}$ with $a_1=1$ and $a_{i+1} =$ the first $\phi$-closed $j \in [a_i+1,n]$
%such that $\DomCard(S_{a_{i+1}} \ldots S_n) > \DomCard(S_{a_i} \ldots S_n)$. That is, if
%$a \in ]a_i, a_{i+1}[$ and $b$ is a dropping index of $a$ then $b \in ]a_i,a_{i+1}[$.

$\Raising(\vec{S}) $ is a finite information, consisting of at most $\MaxInd+1$ relations,
and this no matter what is the number of elements of $\vec{S}$.
$\Raising(\vec{S}) $ can be defined by induction on $\vec{S}$.

%15:39 29/01/2026

In order to inductively define $\Raising(\vec{S}) $, we have to describe how the 
(constant domain-cardinality) 
components of $\minStable(\vec{S})$ change as $\vec{S}$ increases.



\begin{lemma}[Constant Domain-Cardinality Components and Domain-Stability]
\label{lemma-cardinality-stability}
Assume $\vec{S}, S_{n+1}$ is any sequence of injective relations. Suppose 
the components of $\minStable(\vec{S})$ and 
$\minStable(\vec{S}, S_{n+1})$ are, respectively, 
$$
\vec{C_1}, \ldots,\vec{C_k} \ \ \ \mbox{ and } \ \ \   \vec{D_1}, \ldots,\vec{D_h}
$$
with $\vec{C_i} = C_{i,1}, \ldots, C_{i,n_i}$ for $i \in [1,k]$. 
Let us abbreviate for all $i \in [1,k]$: 
$c_i = \DomCard(\vec{C_i})$ and $d_i = \DomCard(\vec{D_i})$.
Let $R$ be the last relation in $\minStable(\vec{S},S_{n+1})$ 
(i.e, the stable tail of $\vec{S},S_{n+1}$).
Then exactly one of the following holds.
\begin{enumerate}
\item
\label{lemma-cardinality-stability-01}
Either $h=k+1$ and
$$
(\vec{C_j} = \vec{D_j} \mbox{ for all $j \in [1,k]$ })  \ \ \  
\mbox{ and } (\vec{D_h} = R)
$$

\item
\label{lemma-cardinality-stability-02}
or for some $i \in [1,k]$ we have
$$
(\vec{C_j} = \vec{D_j} \mbox{ for all $j \in [1,i-1]$}) \ \ \ 
\mbox{ and }     (h=i)  \ \ \ 
\mbox{ and }    (c_{i} > d_{i})    \ \ \ 
\mbox{ and }    (\vec{D_{i}} = R)
$$

\item
\label{lemma-cardinality-stability-03}
or for some $i \in [1,k]$ we have
$$
(\vec{C_j} = \vec{D_j} \ \ \ \mbox{ for all $j \in [1,i-2]$}) \ \ \ 
\mbox{ and }   (h=i-1)  \ \ \ 
\mbox{ and }   (c_i = d_i) \ \ \ 
\mbox{ and }   (\vec{D_{i-1}} = \vec{C_{i-1}}, R)
$$
\end{enumerate}
\end{lemma}

%09:15 05/02/2026

\begin{proof}
We have to prove that $\minStable(\vec{S})$ falls in one of cases 
\ref{lemma-cardinality-stability-01},
\ref{lemma-cardinality-stability-02},
\ref{lemma-cardinality-stability-03}.
above. The components of 
$\minStable(\vec{S}) = \vec{T} = T_1, \ldots, T_m$
are $\vec{C_1}, \ldots,\vec{C_k}$. 
By Lemma \ref{lemma-injective-stable}.\ref{lemma-injective-stable-01} 
we have $c_1 \le \ldots \le c_k$.
Either $\minStable(\vec{S}),S_{n+1}$ is stable or it is not. 

\begin{enumerate}

\item
\emph{Assume $\minStable(\vec{S}),S_{n+1}$ is stable}. Then 
$\minStable(\vec{S},S_{n+1}) = \minStable(\vec{S}),S_{n+1}$ and the
stable tail of $\vec{S},S_{n+1}$ is $S_{n+1}$.
By Lemma \ref{lemma-injective-stable-sequences}.\ref{lemma-injective-stable-sequences-01},
injectivity and stability we have $\DomCard(C_{k,1}) \le \DomCard(S_{n+1})$.

\begin{enumerate}
\item
\emph{Assume $\DomCard(C_{k,1}) < \DomCard(S_{n+1})$}. Then the
components of $\minStable(\vec{S},S_{n+1})$ are in number of $h=k+1$ and
$\vec{D_1}=\vec{C_1}$, \ldots, $\vec{D_k} = \vec{C_k}$ and 
$\vec{D_{k+1}} = S_{n+1}$. This is case \ref{lemma-cardinality-stability-01}.
\item
\emph{Assume $\DomCard(C_{k,1}) = \DomCard(S_{n+1})$}. Then the
components of $\minStable(\vec{S},S_{n+1})$ are in number of $h=k$ and
$\vec{D_1}=\vec{C_1}$, \ldots, $\vec{D_{k-1}} = \vec{C_{k-1}}$ and
$\vec{D_k} = \vec{C_k}, S_{n+1}$. 
This is case \ref{lemma-cardinality-stability-03}.
\end{enumerate}

\item
\emph{Assume $\minStable(\vec{S}), S_{n+1}$ is not stable}. Then 
$\minStable(\vec{S},S_{n+1}) = T_1, \ldots T_j, (T_{j+1} \ldots T_m S_{n+1})$
for the first $j \in [1,m-1]$ such that 
$\DomCard(T_j) > \DomCard(T_j \ldots T_m S_{n+1})$. 
$j$ is one of the raising indexes $a_1, \ldots a_h$ of $\vec{T}$: we have $c_{j-1} \le c_j$
and if it were $c_{j-1} = c_j$, then we would have
$\DomCard(T_{j-1}) = \DomCard(T_j) > \DomCard(T_j \ldots T_m S_{n+1})$,
contradicting the choice of $j$. So $j =a_i$ for some $i \in [1,h]$. 
Let $S=$ the product of $\vec{C_{i}}, \ldots,\vec{C_k}, S_{n+1})$.
%If $i>1$ 
%by Lemma \ref{lemma-injective-stable-sequences}.\ref{lemma-injective-stable-sequences-01}
%we have $\DomCard(C_{i,1}) \le \DomCard(S_{n+1})$.

\begin{enumerate}
\item
\emph{Assume $i=1$ or $\DomCard(S_{a_{i-1}}) < \DomCard(S_{a_i})$}.
Then $\vec{D_j} = \vec{C_j}$ for all $j \in [1,i-1]$ and
$\vec{D_{i}}= R =$ the product of $\vec{C_{i}}, \ldots,\vec{C_k}, S_{n+1}$.
This is case \ref{lemma-cardinality-stability-02}.
\item.
\emph{Assume $i>1$ and $\DomCard(S_{a_{i-1}}) = \DomCard(S_{a_i})$}.
Then $\vec{D_j} = \vec{C_j}$ for all $j \in [1,i-2]$ and 
$\vec{D_{i-1}}= \vec{C_{i-1}}, R$.
This is case \ref{lemma-cardinality-stability-03}.
\end{enumerate}

\end{enumerate}

\end{proof}



%Lemma \ref{lemma-cardinality-stability} above has the following role in the paper. 
%Suppose we have
%\begin{enumerate}
%\item
%a set $\setS$ of sequences of injective relations and a set $\setT$ of  
%raising suffixes of sequences of $\setS$
%\item 
%a measure $\mu:\setT \rightarrow \N$, which is 
%strictly decreasing w.r.t. $\DomCard(\cdot)$. and w.r.t. right-composition.  
%\end{enumerate}
%
%Then if $\Raising(\vec{S})=R_1, \ldots, R_k$ we take as measure $\mu(\vec{S})$ the list 
%$$
%\mu(R_1), \ \ \ \ldots, \ \ \ \mu(R_k)
%$$ 
%of measures of raising suffixes. Lemma \ref{lemma-cardinality-stability} says that that this 
%measure is strictly decreasing w.r.t. $<$. Whenever we add a relation to $\vec{S}$,
%the either we concatenate one more element to the list $\mu(R_1),  \ldots, \mu(R_k)$,
%or the first $R_i$ which changes is replaced with a relation $R$ having a smaller domain-cardinality, hence $\mu(R) < \mu(R_i)$,
%or the first $R_{i-1}$ which changes is replaced with some $R_{i-1} R$ with 
%$\mu(R_{i-1} R) < \mu(R_{i-1})$.


\section{Constant Domain-Cardinality Sequences of Injective Bicolored Relations}

%If $R= (\Stat,Prog)$ is a bicolored sequence, we define $\universe{R} = \Stat \cup \Prog$
%and if $\vec{R}=R_1, \ldots, R_n$ is a bicolored sequence, we define
% $\universe{\vec{R}} = \universe{R_1}, \ldots, \universe{R_n}$.

We have a notion of composition between bicolored relations satisfying $\universe{R S} =
\universe{R} \cup \universe{S}$.

We say that  a bicolored sequence $\vec{R}$ satisfies the injective condition, the domain-stable
condition, the constant domain condition if and only if, respectively, $\universe{\vec{R}}$
satisfies the injective condition, the domain-stable
condition, the constant domain condition. Namely, we move the properties of relation to 
colored relation by forgetting colors.

\begin{definition}[Necessarily Progressing]
Assume $R$ is an colored relation.
We say that $R=(\Stat,\Prog)$ is \emph{necessarily progressing} on $x \in \dom(R)$
if and only if for all $y \in R(x)$ we have $(x,y) \in \Prog$.
\end{definition}

%A constant domain-cardinality bicolored sequence $\vec{R}=R_1, ldots, R_n$
%of injective relations is stationary if and only if for all $i,j \in [1,n]$, $i < j$ we have
%$R_i \ldots R_{j-1} \cap (\dom(R_i) \times \dom(R_j))$ stationary (and 
%a stationary bijection 
%by Lemma \ref{lemma-injective-stable-sequences}.\ref{lemma-injective-stable-sequences-02}).
%
%
%\begin{lemma}[Constant-domain sequence of injective relations]
%\label{lemma-constant-injective}
%A stationary constant-domain bicolored sequence $\vec{C}$ of injective bijections 
%with associated cardinality $\DomCard(\vec{C})=c$ has number of elements 
%$\le \newton{\MaxInd}{c}$.
%\end{lemma}

%%%%%%%%%%%%%%%%%%%%
% NOT YET  CORRECT
%%%%%%%%%%%%%%%%%%%%
%\begin{definition}[Degree of a constant-domain stable colored sequence]
%Assume $\setR$ is a set of injective colored relations on $[1,\MaxInd]$ and
%$R,S$ is a constant-domain $\setR^+$-sequence. 
%
%We call any $d \in \N$ a degree of the pair $R,S$ 
%if and only if for some constant-domain $\setR^+$-sequence 
%$R_1, \ldots, R_n, S$ of relations in $\setR^+$ we have:
%
%\begin{enumerate}
%
%\item
%$R_1, \ldots, R_n, S$ is a constant-domain $\setR^+$-sequence with
%a \emph{stationary} suffix of $d$ elements. 
%
%\item
%$R_1 \ldots R_n = R$.
%
%\end{enumerate}
%\end{definition}
%%%%%%%%%%%%%%%%%%%%%%%%%%%

%17:56 05/02/2026


\begin{definition}[Stationary degree in $\setR^+$]
Assume $\vec{C}= S_1, \ldots, S_e, R_1, \ldots, R_d$ is some 
\emph{non-empty} constant domain-cardinality sequence in $\setR^+$.
Let $d \in \N$. We say that $\vec{C}$ has stationary degree $d$ wr.t. $\setR^+$
if and only if  and:

\begin{enumerate}
\item
$R_i$ is a stationary bijection from $\dom(R_i)$ to $\dom(R_{i+1})$, for all $i \in [1,d-1]$.

\item
no $x \in \dom(R_n)$ is necessarily progressing  
(i.e., for all $x \in \dom(R_n)$ there is some $y$ such that $R_n$ is stationary on $(x,y)$).
\end{enumerate}

We say that $R \in \setR^+$ has stationary degree $d$ wr.t. $\setR^+$
if and only if there is some costant domain-degree sequence 
$\vec{C}= S_1, \ldots, S_e, R_1, \ldots, R_d$ with stationary degree $d$ wr.t. $\setR^+$ 
and such that $R =  S_1 \ldots S_e R_1 \ldots R_d$, namely, $R$ is the product of $\vec{C}$
\end{definition}

We explain the role of stationary degree. 
If $\vec{C}= S_1, \ldots, S_e, R_1, \ldots, R_{d-1}, R_d$ 
has stationary degree $d$ wr.t. $\setR^+$, then in the worst case \emph{all} infinite path $\pi$
passing through $\vec{C}$ could have $d$ consecutive stationary steps.
Indeed, all infinite $\pi$ must move from $\dom(R_i)$ to $\dom(R_{i+1})$ through $R_i$, 
therefore must have $d-1$ consecutive stationary steps. 
Then, no matter what is the step $x \in \dom(R_d)$ of $\pi$, all infinite $\pi$ could have one 
more stationary step, because there is some $y$ with $R_d$ progressing on $(x,y)$.

Let $c = \DomCard(S_1)$ be the domain-cardinality of the last constant domain-cardinality
sequence in an infinite path $\pi$. From $S_1$ on there are $c$ infinite traces in $\pi$.
We will prove that the stationary degree is at most $\newton{\MaxInd}{c}$, therefore
every $\newton{\MaxInd}{c}$ steps some trace progresses. After $n \newton{\MaxInd}{c}$
steps there are at least $n$ progressions, and the maximum number of progressions 
in the $c$ traces is at least $n/c$.
Therefore after infinitely many steps the maximum is infinite and there is some infinitely 
progressing trace.

All $R \in \setR^+$ have at least stationary degree $0$ in $\setR^+$: we take $\vec{C}=S_1$.

Suppose $R  \in \setR^+$ has stationary degree $d$, namely $R$ decomposes as 
$R =  S_1 \ldots S_e R_1 \ldots R_d$ 
with all $R_i \cap \dom(R_i) \times \dom(R_{i+1})$ stationary.  
Then for all $i, j \in [1,d-1]$, $i \le j$ the relation 
$R_i \ldots R_j \cap \dom(R_i) \times \dom(R_{j+1})$
is equal to the product of all 
$
R_i \cap \dom(R_i) \times \dom(R_{i+1})$, \ldots, $R_j \cap \dom(R_j) \times \dom(R_{j+1})
$, 
and therefore it is a stationary bijection.




\begin{lemma}[Degree of a constant-domain sequence of injective colored relations]
\label{lemma-degree-constant}

Assume $\vec{C}= S_1, \ldots, S_e, R_1, \ldots, R_{d}$ is not empty and $R$ has stationary 
degree $d$ in $\vec{C}, \setR^+$ and $c=\DomCard(R)$. 
Assume $\vec{C}, T$ is a constant domain-cardinality sequence in $\setR^+$.

\begin{enumerate}
\item
\label{lemma-degree-constant-01}
Either $R_{d+1}$ is a progressing bijection from $\dom(R_{d+1})$ to $\dom(T)$
or $T$ has some necessarily progressing $x \in \dom(T)$ or $R T$ has degree $d+1$.

\item
\label{lemma-degree-constant-02}
Suppose $\setR^+$ satisfies PDC.
If $R \in \setR^+$ has degree $d$, then $d \le \newton{\MaxInd}{c}$.
\end{enumerate}
\end{lemma}


\begin{proof}
\begin{enumerate}
\item
Suppose $\vec{C}, T$ are constant domain-cardinality sequence, 
and $R = S_1 \ldots S_e R_1 \ldots R_{d} T$.
By Lemma \ref{lemma-injective-stable-sequences}.\ref{lemma-injective-stable-sequences-02}
$R_{d}$ is a bijection from $\dom(R_{d})$ to $\dom(T)$.
If $d=0$ then $R$ has degree $1$. Suppose $d \ge 1$.
Suppose $R_{d}$ is not a progressing bijection from $\dom(R_{d+1})$ to $\dom(T)$;
then it is stationary. Suppose $T$ has no necessarily progressing $x \in \dom(T)$. 
Since all $R_i$ with $i \in [2,d]$ are stationary from $\dom(R_i)$
to $\dom(R_{i+1})$, then $R$ has degree $d+1$.

%10:13 10/02/2026

\item
Suppose $\setR^+$ satisfies PDC and $R \in \setR^+$ has degree $d$. 
Assume $i,j \in [2,d]$, $i < j$.
By assumption, $R_i$ is a stationary bijection from $\dom(R_i)$ to $\dom(R_{i+1})$.
$R_1, \ldots, R_{d+1}$ is constant domain-cardinality sequence, therefore it is stable
and $\dom(R_i \ldots R_{j-1}) = \dom(R_i)$. 

If $\setR^+$ satisfies PDC, then $R_i \ldots R_{j-1}$ is progressing from 
$\dom(R_i \ldots R_{j-1})$ to $\dom(R_i \ldots R_{j-1})$, namely from  $\dom(R_i)$ to 
$\dom(R_i)$. By assumption, $R_i \ldots R_{j-1}$ is \emph{stationary} from  $\dom(R_i)$ to
$\dom(R_{j})$. This forces $\dom(R_i) \not = \dom(R_{j})$, for all $i,j \in [2,d]$, $i < j$.
Let $c = \DomCard(R_1)$ the common domain cardinality of all $R_2, \ldots, R_{d+1}$. 
There are $\newton{\MaxInd}{c}$ subsets of $[1,\MaxInd]$
of cardinality $c$, and $d$ different domains for $R_2, \ldots, R_{d+1}$,
therefore $d \le \newton{\MaxInd}{c}$. 
\end{enumerate}
\end{proof}

In the proof above, if $d$ is any of the degrees of
$\vec{C}$, then $(\newton{\MaxInd}{c}-d) \in [0,\newton{\MaxInd}{c}]$.
We can extend $\vec{C}$ with stationary bijections to 
a constant domain-cardinality sequence for at most
$\newton{\MaxInd}{c}-d$ times. When $\newton{\MaxInd}{c}-d=0$, in 
all constant domain-cardinality extensions $R_1, \ldots, R_{d+1},T$
the relation $R_{d+1}$ is progressing between $\dom(R_{d+1})$ and $\dom(T)$.


%Assume $\vec{S} = S_1, \ldots, S_n$ is any sequence of bicolored relations. We define
%the uncolored version of $\vec{S}$ as $\universe{\vec{S}} = \universe{S_1}, \ldots, 
%\universe{S_n}$ and the minimum domain-stable grouping of $\vec{S}$ as
%$\minStable(\vec{S}) = \minStable(\universe{\vec{S}})$.
%
%We define the integer sequence $\sigma(i) = \Card(S_i \ldots S_n)$ for $i=1, \ldots, n$.
%If $S_1, \ldots, S_n$ are injective,  then $\sigma$ is weakly increasing.
%Suppose $\sigma([1,n]) = \{c_1, \ldots, c_h\}$ for some increasing sequence 
%$0 \le c_1 < c_2 < \ldots < c_h$ (i.e., if we skip repeated values). 
%For $i=1, \ldots, h$, let $a_i$ be the last index such that
%$ \sigma(a_i) = c_i$. If $S_1, \ldots, S_n$ are injective, then $a_i$ is increasing
%and $\sigma$ decomposes in $h$ adiacent constant subsequences:
%$$
%\sigma(1), \ldots, \sigma(a_1),
%\ \ \
%\sigma(a_1+1), \ldots, \sigma(a_2),
%\ \ \ 
%\ldots,
%\ \ \
%\sigma(a_{h-1}+1), \ldots, \sigma(a_h)
%$$ 
%The subsequence number $i$ has constant value $c_i$.

%Suppose:
%\begin{enumerate}
%\item
%$R_1, \ldots, R_n \in \setR^+$ (hence each $R_i$ is injective and satisfies PDC) 
%\item
%Suppose $R_1, \ldots, R_n$ is stable. 
%\end{enumerate}
%Set $D_i = \dom(R_i)$ for all $i \in [1,n]$.


%%%%%%%%%%%%%%%%%%%%%%%%%%
% 20:36 24/01/2026
%%%%%%%%%%%%%%%%%%%%%%%%%%
% and $D_n = \codom(R_n)$ if $R_n$ is a bijection 
% and $D_n = \dom(R_n)$ otherwise.
%%%%%%%%%%%%%%%%%%%%%%%%%%
% (i). forse è meglio mettere il caso particolare per D_n tra i possibili miglioramenti
% per la leggibilità della prova ottenuta, anziché nelle definizione di base
%%%%%%%%%%%%%%%%%%%%%%%%%%
% (ii). Altro miglioramento: se il massimo valore del grado \deg(\vec{R},j) per j 
% ultima posizione della componente numero h è M allora si può rimpiazzare 
% 2^\MaxInd e (2^\MaxInd - 1 - \deg(\vec{R},j))  
% con M+1 e con (M - \deg(\vec{R},j)) 
%%%%%%%%%%%%%%%%%%%%%%%%%%%
% (iii). Altro miglioramento: se un elemento della componente di cardinalità h ha
% una sola immagine nella componente successiva, si può eliminare nella
% nella componente successiva, perché se decresce entro la componente di cardinalità h
% allora la componente di cardinalità h decresce e le successive sono irrilevanti.
%%%%%%%%%%%%%%%%%%%%%%%%%%%


%\begin{lemma}
%\begin{enumerate}
%\item
%If $i,j \in [1,n]$, $i<j$ and $D_i = D_j$ then $R_i \ldots R_j$ is progressing on $D_i$.
%\item
%If  $i,j \in [1,n]$, $i<j$ and $\Card(D_i) = \Card(D_j)$ and $j-i \ge 2^{\MaxInd}$ 
%then $R_i \ldots R_j$ is progressing on $D_i$.
%\end{enumerate}
%\end{lemma}


\section{The Size Map}

We assume having a size map $\size : ProofNodes, 
[1,\MaxInd] \rightarrow \N^+$, returning an expression $\sigma(\alpha,i)$ denoting a 
\emph{positive} integer. 
Assume $p(\vec{\alpha})$ is the inductive formula number $j$ of node $\alpha$ and
$q(\vec{u})$ is the inductive formula number $j$ of node $\beta$.
Assume $\beta$ is the father of $\alpha$ and $\alpha$ the conclusion of an inductive rule.

We write $\size(\beta,j) < \size(\alpha,j)$ for: for all closed substitution $\sigma$ we have
 $\sigma(\size(\beta,j)) < \sigma(\size(\beta,j)) $. 
 
We require that: if $q(\vec{u})$
an assumption of the inductive rule of conclusion
$p(\vec{t})$ then $\size(\beta,j) < \size(\alpha,j)$. 

We often skip $\alpha$ and we just write $\size(i), \size(j)$.


We include some examples of size map, counting the
inductive rules in the proof of a closure of $q(\vec{t})$.

\begin{enumerate}
\item 
if $q(\vec{t}) = \N(v)$, then $$\size(i)=v+1$$ If $v$ has value $a$ 
then $v=S^a(0)$ and the proof of $\N(v)$ has $a$ times an $S$-rule and $1$-time
the $0$-rule.

\item
If $q(\vec{t}) = \List(l)$ states that $l$ is a list on $\N$, 
then $$\size(i) = (\length(l)+1)+\Sigma_{i=1}^{i=\length(l)}(l[i]+1)$$
Here we count $\length(l)$ times the rule $\cons$, once the rule $\nil$,
the we add the total size of all elements of $l$.

\item
If $q(\vec{t}) = \Tree(t)$ states that $t$ is a binary tree with labels in $\N$, 
then $\size(i)$ is any expression counting the nodes in $t$, plus the size
of all labels of the nodes $t$.
\end{enumerate}

Alternatively, the size map can count the maximum length of sequence of 
inductive rules in the proof of $q(\vec{t})$. In this case the size of $\N(v)$ is $v+1$ again,
but the size of a list or of a tree can be smaller.
\\


\section{Size of a Set and of a Relation}
Assume $\vec{S}$ has minimum stable grouping 
$\minStable{\vec{C}} = \vec{R} = R_1, \ldots, R_n$. Assume there are $h$ maximal
constant-cardinality domain components, and the component n umber $i$ is  
$R_{a_i}, \ldots ,R_{b_i}$, with $c_i  = \Card(R_{a_i})$.  


We extend the size map to sets of integers in $[1,\MaxInd]$,
These integers are position inductive formulas in the left-hand side of a node $\alpha$,
We extend the size map to relations $R$ associated to a path between two nodes 
$\alpha$ and $\beta$.
 
 %16:19 06/02/2026

\begin{definition}[Size of Sets and Relations]
Assume $X \subseteq [1,\MaxInd]$ and $R:\alpha \rightarrow \beta$  is a colored
relation.
\begin{enumerate} 
\item
We set $\size(\alpha, X) = \Sigma_{i \in X} \size(\alpha,i)$.
\item
For each relation $S$ recall that $S(x) = \{y \in [1,\MaxInd] | x S y \}$. 
If $R$is colored we define $R(x) = \universe{R}(x)$, forgetting colors.
\item
We associate to each (colored or uncolored) $R:\alpha \rightarrow \beta$ 
an integer expression $\size(\beta, R) =
\Sigma_{x \in \dom(R)} \max( \{ \size(\beta,y) | y \in R(x) \} )$.
\end{enumerate}
\end{definition}

\begin{lemma}[Decreasing Size]
\label{lemma-decreasing-size}
Assume $R:\alpha \rightarrow \beta$, $S:\beta \rightarrow \gamma$ 
are colored relations.
\begin{enumerate}
\item
$\size(\gamma, RS) \le \size(\beta, R)$.
\item
If $\dom(RS) \subset \dom(R)$ then
$\size(\gamma, RS) < \size(\beta, R)$.
\end{enumerate}
\end{lemma}

\begin{proof}
\begin{enumerate}
\item
We have $\size(\gamma, RS) =
\Sigma_{x \in \dom(RS)} \max(\size(\gamma, RS(x))$. For all $x \in \dom(RS)$ we have
$x \in \dom(R)$ and if $z \in RS(x)$ then for some $y \in R(x)$ we have $y S z$,
therefore $\size(\beta,y) \ge \size(\gamma,z)$. We deduce 
$\max( \{ \size(\beta,z) | z \in RS(x) \}\le \max( \{ \size(\beta,y) | y \in R(x) \}$, 
hence 
$
\size(\gamma, RS)                                                                              =
\Sigma_{x \in \dom(RS)} \max( \{ \size(\beta,z) | z \in RS(x) \})     \le
\Sigma_{x \in \dom(RS)} \max( \{ \size(\beta,y) | y \in R(x) \}        \le
$
(by $\dom(RS) \subseteq \dom(R)$)
$
\Sigma_{x \in \dom(R)}   \max( \{ \size(\beta,y) | y \in R(x) \}        =
\size(\beta, R)
$.
\item
Assume $\dom(RS) \subset \dom(R)$.
Then 
$
\Sigma_{x \in \dom(RS)} \max( \{ \size(\beta,y) | y \in R(x) \}       <
\Sigma_{x \in \dom(R)}   \max( \{ \size(\beta,y) | y \in R(x) \}
$, 
because the second sum has more addends. We conclude that 
$
\size(\gamma, RS)                                                                              =
\Sigma_{x \in \dom(RS)} \max( \{ \size(\beta,z) | z \in RS(x) \})     \le
\Sigma_{x \in \dom(RS)} \max( \{ \size(\beta,y) | y \in R(x) \}        <
\Sigma_{x \in \dom(R)}   \max( \{ \size(\beta,y) | y \in R(x) \}        =
\size(\beta, R)
$, as wished.
\end{enumerate}
\end{proof}

If $\vec{S} = S_1, \ldots, S_n \in \setR^+$ is stable then we will define a notion of
size for $\vec{S}$.
For a generic injective sequence $\vec{R}$, instead, we take the stable version 
$\vec{S}$ of $\vec{R}$ first.


\begin{definition}[Size of an Injective Sequence]
Suppose $\dom(S)$ is any injective domain-stable sequence in $\setR^+$,
with $h$ components $\vec{C_i} = S_{a_i},\ldots, S_{b_i}$ for $i \in [1,h]$. 
Suppose the product $R_i = S_{a_i}\ldots S_{b_i}$ of $\vec{C}$
has domain-cardinality $c_i$ and degree $d_i$ w.r.t. $\vec{C_i}$, $\setR^+$. Then the list 
$$
(
( \size(S_{a_1}\ldots S_n), \newton{\MaxInd}{c_1}-d_1 ), 
\ldots,
( \size(S_{a_h}\ldots S_n), \newton{\MaxInd}{c_h}-d_h )
)
$$ 
is a size of $\vec{S}$.

If $\dom(R)$ is any injective sequence then any size of 
$ \size(\minStable(\vec{R}))$ is a size of $\vec{R}$.
\end{definition}


\begin{lemma}[Length Bound from PDC, IC, Constant-Cardinality Domain]
\label{lemma-length-bound}
\begin{enumerate}
\item
\label{lemma-length-bound-01}
For all constant-cardinality domain colored sequence $R_1 \ldots R_n R_{n+1}$ we have
$\size(R_1 \ldots R_n R_{n+1}) = \size(R_{n+1})$.

\item
\label{lemma-length-bound-02}
For all colored relations $R:\alpha \rightarrow \beta$: $size(\alpha,\dom(R)) \ge \size(\beta,R)$.
 
\item
\label{lemma-length-bound-03}
Assume $R_1 \cap (\dom(R_1) \times \dom(R_2))$ is a bijection.
Then $\size(R_1) \ge \size(\dom(R_2)) \ge \size(\beta,R)$.

\item
\label{lemma-length-bound-04}
Assume $R_1,R_2$ is a constant-cardinality domain sequence 
of injective colored relations.
If $\size(R_1) = \size(R_2)$, then $\size(R_1) = \size(\dom(R_2)) = \size(R_2)$
% and $R_1 \cap (\dom(R_1) \times \dom(R_2))$ is a stationary bijection.

\item
\label{lemma-length-bound-05}
Assume $R_1, \ldots, R_n, R_{n+1}$ is a constant domain-cardinality sequence of 
relations in $\setR^+$ with domain-cardinality $c$.
Suppose for some $k$
we have $\size(R_i) = k$ for all $i \in [1,n+1]$. 

Then $\size(\dom(R_i)) = k$ for all $i \in [2,n+1]$ and  
$R_i \cap (\dom(R_i) \times \dom(R_{i+1}))$ is a stationary bijection for all $i \in [2,n]$ and
$$
n \le \newton{\MaxInd}{c}
$$ 
\end{enumerate}
\end{lemma}


\begin{proof}

\begin{enumerate}
\item
%\label{lemma-length-bound-01}
We assume that $R_1, \ldots, R_n, R_{n+1}$ is some colored constant-cardinality domain 
sequence  with $R = R_1 \ldots R_n R_{n+1}$ in order to prove that
$\size(R) = \size(R_{n+1})$.
By Lemma \ref{lemma-domain-stability}.\ref{lemma-domain-stability-02}, 
$\dom(R) = \dom(R_1)$.
By Lemma \ref{lemma-injective-sequences}.\ref{lemma-injective-sequences-01},
$R_1 \ldots R_n$ restricted to $\dom(R_1)$ and  $\dom(R_{n+1})$ 
is the graph of some bijection $\psi$. Therefore for
all $x \in \dom(R_1)$ we have $R(x) = R_1 \ldots R_n R_{n+1}(x) = R_{n+1}(\psi(x))$.

If $\psi:X \rightarrow Y$ is a bijection, for any sum $\Sigma_{x \in X} a_{\psi(x)}$, 
we have $\Sigma_{x \in X} a_{\psi(x)} = \Sigma_{y \in Y} a_y$. 

Thus, from $\dom(R) = \dom(R_1)$ and $x R z \equiv \psi(x) R_{n+1} z$ and $\psi$ bijection 
we deduce that 
$$
\size(R) =
\Sigma_{x \in \dom(R)}             \max(\{\size(z) | z \in R(x) \})                            = 
\Sigma_{x \in \dom(R_1)}         \max(\{\size(z) | z \in R_{n+1}(\psi(x) )  \})     =
\Sigma_{y \in \dom(R_{n+1})} \max(\{\size(z) | z \in R_{n+1}(y)           \})      =
\size(R_{n+1})
$$

\item
%\label{lemma-length-bound-02}
For all $x \in \dom(R)$, $y \in R(x)$ from $x R y$ and the conditions 
about coloring, we deduce that $\size(x) \ge size(y)$. 
Thus, $\size(x) \ge \max(\{\size(z) | z \in R(x) \}$.
We conclude that $size(\dom(R)) = \Sigma_{x \in \dom(R)} \size(x) \ge 
\Sigma_{x \in \dom(R)} \max(\{\size(z) | z \in R(x) \}) = \size(R)$.

\item
%\label{lemma-length-bound-03}
If $R_1 \cap (\dom(R_1) \times \dom(R_2))$ is a bijection $\psi$,
then for all $x \in \dom(R_1)$ there is exactly  one $y \in \dom(R_2)$ with $y = \psi(x)$.
Thus, $ \size(R_1) = \Sigma_{x \in \dom(R_1)} \max(\{\size(y) | y \in R_1(x) \}) \ge 
 \Sigma_{x \in \dom(R_1)} \size(\psi(x)) =  
\Sigma_{y \in \dom(R_2)} \size(y) = \size(\dom(R_2))$.

We have $\size(\dom(R_2)) \ge \size(R_2)$ by point \ref{lemma-length-bound-02} above.

\item
%\label{lemma-length-bound-04}
By $R_1, R_2$ constant domain-cardinality sequence and $R_1, R_2$ injective and
Lemma \ref{lemma-injective-sequences}.\ref{lemma-injective-sequences-01} 
we deduce that $R_1 \cap (\dom(R_1) \times \dom(R_2))$ is a bijection.
By point \ref{lemma-length-bound-03} above we deduce that 
$\size(R_1) \ge \size(\dom(R_2)) \ge \size(R_2)$, and from $\size(R_1) = \size(R_2)$ 
we conclude that $\size(R_1) = \size(\dom(R_2)) = \size(R_2)$. 

\item
%\label{lemma-length-bound-05}
By Lemma \ref{lemma-size-colored-relation} $R_i \cap \dom(R_i) \times \dom(R_{i+1})$ 
is a stationary bijection for all $i \in [1,n ]$. By the previous point we have 
$\size(\dom(R_i)) = k$ for all $i \in [2,n+1]$. Thus, $R_i \cap (\dom(R_i) \times 
\dom(R_{i+1}))$ is a stationary bijection for all $i \in [2,n]$, otherwise we would have
$\size(\dom(R_i)) > \size(\dom(R_{i+1}))$.
Lemma \ref{lemma-degree-constant}.\ref{lemma-degree-constant-02}
implies that $n \le \newton{\MaxInd}{c}$.
\end{enumerate}
\end{proof}

\begin{lemma}[Size of a Colored Relation Eventually Decreases]
\label{lemma-size-colored-relation}
Assume that:

\begin{enumerate}
\item
$\setR^+$ satisfies PDC and the Injectivity Condition IC. 
\item 
$R_1, \ldots, R_n, R_{n+1}$ is a constant domain-cardinality sequence of 
relations in $\setR^+$. 
\item
Suppose $\size(R_1 \ldots R_n) = \size(R_1 \ldots R_n R_{n+1})$.
\end{enumerate}

Then for all $i, j \in [1,n]$, $i \le j$:
If $\size(R_1 \ldots R_n) = \size(R_1 \ldots R_n R_{n+1})$,
then $\size(R_n) = \size(\dom(R_n) = \size(R_{n+1})$,
and if $R_1 \ldots R_n$ has stationary degree $d$, then 
$R_1 \ldots R_{n+1} $ has stationary degree $d+1$.

\end{lemma}

\begin{proof}
By Lemma \ref{lemma-length-bound}.\ref{lemma-length-bound-01}
from $\size(R_1 \ldots R_n) = \size(R_1 \ldots R_n R_{n+1})$
we deduce that $\size(R_n) = \size(R_{n+1})$.

By Lemma \ref{lemma-length-bound}.\ref{lemma-length-bound-04}
we conclude that $\size(R_n) = \size(\dom(R_{n+1})) = \size(R_{n+1})$.

Assume $R_1 \ldots R_n$ has stationary degree $d$.
If $d=0$, then $R_1, \ldots, R_{n+1} $ includes at least two relations and therefore it
has stationary degree $1$. Now assume $d \ge 1$. By definition of state, 
$\size(\dom(R_{n})) = \size(R_n)$, therefore 
 $ \size(\dom(R_{n})) = \size(\dom(R_{n+1})) $. Thus,
 $R_n \cap \dom(R_n) \times \dom(R_{n+1})$ is a stationary bijection
and $R_1 \ldots R_{n+1} $ has stationary degree $d+1$.
\end{proof}

%We claim that after at most $2^\MaxInd$ steps any constant-domain sequence progresses.
%We formally express this fact through the notion of degree.
% 
%\begin{definition}[Degree of an Index in a Domain-Stable Injective Sequence]
%The degree $\deg(j,\vec{R})$ of $j \in [1,n]$ in 
%$\vec{R} = R_1, \ldots, R_n$ as above is $2^\MaxInd-1-d$, for $d$ the minimum distance of
%an index $i \le j$ such that at least one of the following happens:
%\begin{enumerate}
%\item
%$i=1$
%\item
%$i>1$ and $\Card(D_{i-1}) < \Card(D_j)$
%\item
%$R_{i} \ldots R_j$ is progressing between $D_i$ and $D_j$.
%\end{enumerate}
%\end{definition}
 

\section{State of a Path}
We define a graph of states and transitions for the algorithm translating a cyclic proof. We first
associate a state to each sequence of injective relations.

%19:11 31/01/2026

\begin{definition}[The Set $\setS$ of States]
A state is any sequence $\sigma \Seq(\setR^+ \times \N \times \setR^+)$ 
on $\setR^+ \times \N \times \setR^+$:
$$
\sigma 
=
(
R_1, d_1, S_1;
\ldots;
R_h, d_h, S_h
)
$$
such that:
\begin{enumerate}
\item
there are some $T_1, \ldots, T_h \in \setR+ \cup \{\id_{[1,\MaxInd]}\}$ such that 
for all $i \in [1,h-1]$:
\begin{enumerate}
\item
if $i<h$ then $R_i = T_i S_i R_{i+1}$
\item
if $i=h$ then $R_h = T_h S_h$.
\item
$d_i$ is a degree of $T_i S_i$ in $\setR^+$.
\end{enumerate}
\item
For all $i,j \in [1,d]$, $i<j$ we have $\DomCard(R_i) < \DomCard(R_j)$.
\end{enumerate}
$\States(\setR^+) = 
\{\sigma \in \Seq(\setR^+ \times \N \times \setR^+)
 | \sigma \mbox{ is a state} \}$ 
 is the set of states of $\setR^+$.
\end{definition}

\begin{definition}[State of an Injective Domain-Stable Sequence in $\setR^+$]
Assume  $\setR^+$ is any set of injective relations and $\vec(S) = S_1, \ldots, S_n$ 
is any domain-stable sequence in $\setR^+$ with 
components $\vec{C_i} = S_{a_i},\ldots, S_{b_i}$ for $i \in [1,h]$.
Suppose $C_i = S_{a_i}\ldots S_{b_i}$
\begin{enumerate}
\item
We say that a sequence $\sigma$ on $\setR^+ \times \N \times \setR^+$ is a state of 
$\vec{S}
= S_1, \ldots, Sn$ 
in $\setR^+$ and we write $\sigma \state \vec{S}$ if and only if for some degree $d_i$ of  
$\vec{C_i}$ w.r.t. $\setR^+$ we have
$$
\sigma \ \ \ 
= \ \ \ 
(
S_{a_1} \ldots S_{n }, d_1, S_{b_1};
\ldots;
S_{a_h} \ldots S_{n }, d_h, S_{b_h}
)
$$
% THIS IS BUT A CONSEQUENCE OF THE DEFINITION OF DEGREE
% and for all $i \in [1,h]$, either $\size(\dom(R_{b_i})) = \size(R_{b_i})$ or $d_i=0$.

%CONTROLLARE 20:55 09/02/2026

\item
If $\vec(R)$ is any injective sequence then $\sigma \state\vec(R)$ in $\setR^+$ if and only
if $\sigma \state \minStable(\vec{R})$ in $\setR^+$. 

\item
If $\sigma \state \minStable(\vec{R})$ in $\setR^+$ then 
$$
\size(\sigma)
=
(
\size(C_1), \newton{\MaxInd}{c_1}-d_1; 
\ldots;
\size(C_h), \newton{\MaxInd}{c_h}-d_h
)
$$
\end{enumerate}
\end{definition}

Remark that we can decide whether $\size(\dom(R)) = \size(R)$, by checking whether
no $x \in \dom(R)$ is necessarily progressing
for all $x \in \dom(R)$ there is some $y \in R(x)$ such that $R$ is stationary on $(x,y)$.
When $\size(\dom(R_{b_i})) = \size(R_{b_i})$ we will reset the counter $d_i$ to $0$.


\begin{lemma}[State of a Sequence]
Assume  $\setR^+$ is any set of injective relations and
$\sigma \state \vec(R)$ in $\setR^+$. Then $\sigma \in \States(\setR^+)$.
\end{lemma}

\begin{proof}
Assume  $\setR^+$ is any set of injective relations and
$\sigma \state \vec(R)$ in $\setR^+$, in order to check the clauses for 
$\sigma \in \States(\setR^+)$.
Suppose $\minStable(\vec{R}) = \vec{S} = S_1,\ldots, S_n$ and
$$
\sigma \ \ \ 
= \ \ \ 
(
S_{a_1} \ldots S_{n }, d_1, S_{b_1};
\ldots;
S_{a_h} \ldots S_{n }, d_h, S_{b_h}
)
$$
We check the definition of state.
\begin{enumerate}
\item
there are some $T_1=S_{a_1} \ldots S_{b_1-1}, \ldots, T_h=S_{a_h} \ldots S_{b_h-1} \in \setR+ \cup \{\id_{[1,\MaxInd]}\}$ such that 
for all $i \in [1,h-1]$:
\begin{enumerate}
\item
if $i<h$ then 
$S_{a_i} \ldots S_{n } = (S_{a_i} \ldots S_{b_i-1}) S_{b_i} (S_{a_{i+1}} \ldots S_{n })$.
\item
if $i=h$ then 
$S_{a_h} \ldots S_{n } = (S_{a_h} \ldots S_{b_h-1}) S_{b_h}$.
\item
$d_i$ is a degree of $(S_{a_i} \ldots S_{b_i-1}) S_{b_i}$ in $\setR^+$ by
definition of $\sigma \state \vec(R)$.
\end{enumerate}
\end{enumerate}
\end{proof}

%By definition unfolding, if $\dom(R) = R_1, \ldots, R_n$ is any injective sequence, 
%$\minStable(\vec{R})$ has components  
%$S_{a_i},\ldots, S_{b_i}$ for $i \in [1,h]$, and $\sigma$
%is a state of $\vec{R}$, then $S_{a_i}\ldots S_{b_h} = R_{a_i}\ldots R_{n}$.
%$S_n$ is $R_k \ldots R_n$ for the first $k \in [1,n]$
%such that $\phi(k)=n$: we can have $k=n$. 
%We deduce that $\size(\sigma)$ is
%
%$$
%(
%\size(R_{a_1}\ldots R_{n}), d_1, 
%\ldots,
%\size(R_{a_h}\ldots R_{n}), d_h,
%R_k \ldots R_n
%)
%$$

%17:00 06/02/2026



%\begin{definition}
%\label{definition-transition}
%Assume $\vec{R}, R$ is any sequence of injective relations in $\setR^+$, 
%with $c = \DomCard(R)$. Suppose a state $\sigma$ of $\vec{R}$ in $\setR^+$ is
%$$
%(
%(\size(R_{a_1}\ldots R_{n}), d_1), 
%\ldots,
%(\size(R_{a_h}\ldots R_{n}), d_h)
%)
%$$ 
% and the constant domain-cardinality components of
%$\minStable(\vec{R}, R)$ are, 
%$$
%\vec{D_1}, \ldots,\vec{D_h}
%$$
%Then either $h=k+1$ and $\trans(\sigma,R)$ is
%
%$$
%(
%(\size(R_{a_1}\ldots R_{n} R), d_1), 
%\ldots,
%(\size(R_{a_h}\ldots R_{n} R), d_h),
%(\size(R),                                     \newton{\MaxInd}{c})
%)
%$$ 
%
%or $k \in [1,h]$ and if 
%$c_i = \DomCard(\vec{C_i}))$, 
%$d_i = \DomCard(\dom(\vec{D_i})$ for all $i \in [1,k]$ 
%and $R_{a_k}\ldots R_{n} <_1 R_{a_k}\ldots R_{n}, R$ 
%is a progressing extension or $(c_k < d_k)$ and $\trans(\sigma,R)$ is
%
%$$
%(
%(\size(R_{a_1}\ldots R_{n} R), d_1), 
%\ldots,
%(\size(R_{a_{k-1}}\ldots R_{n} R), d_{k-1}),
%(\size(R_{a_k}\ldots R_{n} R), \newton{\MaxInd}{c})
%)
%$$ 
%
%or $R_{a_k}\ldots R_{n} <_1 R_{a_k}\ldots R_{n}, R$ is a stationary extension and 
%$(c_k = d_k)$ and
%
%$$
%(
%(\size(R_{a_1}\ldots R_{n} R),         d_1), 
%\ldots,
%(\size(R_{a_{i-1}}\ldots R_{n} R),  d_{k-1}),
%(\size(R_{a_{i}}\ldots R_{n} R),     d_k + 1)
%)
%$$ 
%
%\end{definition}

%\begin{enumerate}
%\item %1
%Suppose$ \vec{R}S$ is stable and $\DomCard(S)> c_h$. Then
%$c_{h+1} = \DomCard(S)$, $a_h = n$, $a_{h+1} = n+1$ and 
%the state of $\vec{S}S$ is 
%
%$$
%(R_{a_1}\ldots R_{b_1}, d_1, 
%\ldots,
%R_{a_h}\ldots R_{b_h}, d_h,
%S, \newton{\MaxInd}{c_{h+1}})
%$$
%
%\item %2
%Suppose$ \vec{R}S$ is stable and $\Card(\dom(R)= c_h$
%and $R_n S$ is progressing on $\dom(R_{n})$.
%Then the new value of $a_h$ is $n+1$ and $\state(\vec{S}S)$ is 
%
%$$
%(c_1, R_{a_1}\ldots R_{a_2-1}, \deg(a_1,\vec{R}), 
%\ldots, 
%c_{h-1}, R_{a_{h-1}}\ldots R_{a_h}, \deg(a_h,\vec{R}),
%c_{h}, S, 2^{\MaxInd} - 1)
%$$
%
%\item %3
%Suppose$ \vec{R}S$ is stable and $\Card(\dom(R)= c_h$
%and $R_n S$ is stationary on $\dom(R_{n})$.
%
%Then $\state(\vec{S}S)$ is 
%
%$$
%(c_1, R_{a_1}\ldots R_{a_2-1}, \deg(a_1,\vec{R}), 
%\ldots,
%c_{h-1}, R_{a_{h-1}}\ldots R_{a_h}, \deg(a_h,\vec{R}),
%c_{h}, S, \deg(a_h,\vec{R}) - 1)
%$$
%
%\item %4
%Suppose$ \vec{R}S$ is not stable and $j \in [1,n]$ is the first index
%such that $\dom(R_j \ldots R_n) \supset \dom(R_j \ldots R_n S)$.
%Then $j = a_i + 1$ for some $i$ and the Minimum Stable grouping of $\vec{S}S$
%is $R_1, \ldots, R_{a_i}, R$ with $R = R_j \ldots R_n S$. If $R_{a_i} R$
%is progressing on $\dom(R_{a_i})$ then the new value of $a_i$ is
%$a_{i+1}$ and $\state(\vec{S}S)$ is 
%
%$$
%(c_1, R_{a_1}\ldots R_{a_2-1}, \deg(a_1,\vec{R}), 
%\ldots,
%c_{i-1}, R_{a_{i-1}}\ldots R_{a_i}, \deg(a_i,\vec{R}),
%c_{i}, R, 2^{\MaxInd}-1)
%$$
%
%\item %5
%If instead $R_{a_i} R$
%is stationary on $\dom(R_{a_i})$ then $\state(\vec{S}S)$ is:
%
%$$
%(c_1, R_{a_1}\ldots R_{a_2-1}, \deg(a_1,\vec{R}), 
%\ldots,
%c_{i-1}, R_{a_{i-1}}\ldots R_{a_i}R, \deg(a_i,\vec{R})-1)
%$$
%
%\end{enumerate}

\subsection{The Graph of States of the Translation Algorithm}
Assume be given a set $\setR$ of colored relations on $[1,\MaxInd]$. We assume
that $\setR^+$ satisfies the Injectivity Condition IC and the Progressing Domain Condition
PDC. 

A preliminary step for the translation algorithm is to derive the closure $\setR^+$ of $\setR$ 
by compatible composition: $\setR^+$ is the smallest set of relations such that 
$\setR \subseteq \setR^+$ and for all $R:\alpha \rightarrow \beta \in \setR^+$ and 
$S : \beta \rightarrow \gamma \in \setR^+$ we have 
$RS:\beta \rightarrow \gamma \in \setR^+$.

$\setR^+$ can be computed in time $2^{O(\MaxInd^2)}$, more precisely with 
$2^{3*\MaxInd^2}$ operations. The rest of the algorithm requires time $2^{O(\MaxInd^3)}$,
hence exponential time.
\\

%12:58 08/02/2026

%\begin{definition}[The Translation Algorithm State]
%A state $\sigma$ of the Translation algorithm is either the empty sequence $()$,
%or any sequence $(R_1, d_1, \ldots, R_n, d_n; S)$ for some $n \in \N$, 
%$R_1, \ldots, R_n, S \in \setR^+$ and $d_1, \ldots, d_n \in \N$, and such that:
%
%\begin{enumerate}
%\item
%for all $i \in [1,n-1]$, $\DomCard(R_i) \not = \DomCard(R_{i+1})$ 
%and there is some $S_i \in \setR^+$ such that $R_i = S_i R_{i+1}$
%
%\item
%$\DomCard(R_n) = \DomCard(S)$ and for all $i \in [1,n]$ there is some $T_i \in \setR^+$ 
%such that $R_i = T_i S$
%\end{enumerate}
%
%$()$ is called the \emph{initial state}.
%\end{definition}

%By Lemma \ref{lemma-injective-sequences}.\ref{lemma-injective-sequences-01}, 
%injectivity imply that $\DomCard(S_i R_{i+1}) \le \DomCard(R_{i+1})$, 
%therefore if $R_i = S_i R_{i+1}$ then we have $\DomCard(R_i) \le \DomCard(R_{i+1})$. 
%Thus, the condition $\DomCard(R_i) \not = \DomCard(R_{i+1})$ is equivalent to
%$\DomCard(R_i) < \DomCard(R_{i+1})$.

Suppose we add any $R : \beta \rightarrow \gamma \in \setR^+$ 
to the sequence $\vec{R}:\alpha \rightarrow \beta$ in $\setR^+$.
For each state $\sigma \in \States(\setR^+)$ we define a state $\tau = \trans(\sigma,R) \in \States(\setR^+)$
such that if $\sigma \state \vec{R}$ then $\tau \state \vec{R},R$ and $\sigma > \tau$ in the lexicographic ordering.

%In the next definition we consider the same cases of Lemma \ref{lemma-cardinality-stability}. 


\begin{definition}[The Transition Function $\trans(\sigma,R) \in \States(\setR^+)$]
\label{definition-transition-state}

Suppose $\sigma = (U_1, d_1,V_1; \ldots; U_h, d_h V_h) \in \States(\setR^+)$ and
$R \in \setR^+$. We set
$$
U'_i = U_i R \mbox{ for all $i \in [1,h]$ \ \ \ and } U'_{h+1} = R
$$ 
Suppose that $i$ is the first element of $i \in [1,h+1]$ such that:

\begin{enumerate} 
\item
either $i \in [1,h]$ and $\dom(U'_i) \subset \dom(U_i)$
\item 
or $i=h+1$
\end{enumerate}

We define a transition function $\trans(\sigma,R)$. 

\begin{enumerate}

\item 
\emph{Assume that $i=1$ or $\DomCard(U'_i) \not = \DomCard(U'_{i-1})$}.
Then $\trans$ sets the counter number $i$ to its minimum value $0$

$$
\trans(\sigma,S) = (U'_1, d_1, V_1; \ \ldots; U'_{i-1}, d_{i-1},V_{i-1}; \ U'_i, 0, \ U'_i)
$$

\item
\emph{Assume that $i>1$ and $\DomCard(U'_i) = \DomCard(U'_{i-1})$}.
\begin{enumerate}
\item
\emph{Assume $\size(\dom(U'_{i})) \not = \size(U'_i)$ or 
($d_{i-1}\ge 1$ and $V_{i-1}$ is progressing from $\dom(V_{i-1})$ to 
$\dom(U'_{i})$)}. Then $\trans$ sets the counter number $(i-1)$ to $0$

$$
\trans(\sigma,S) = 
(U'_1, d_1,V_1;  \ \ldots; U'_{i-2}, d_{i-1},V_{i-2};  \ U'_{i-1}, 0, U'_{i-1})
$$

\item
\emph{Assume $\size(\dom(U'_{i})) = \size(U'_i)$ and
($d-{i-1}= 0$ or $V_{i-1}$ is stationary from $\dom(V_{i-1})$ to $\dom(U'_{i})$)}.
Then $\trans$ increases the counter number $(i-1)$ by one:

$$
\trans(\sigma,S) = 
(U'_1, d_1,V_1;  \ \ldots; U'_{i-2}, d_{i-1},V_{i-2};  \ U'_{i-1}, d_{i-1}+1, U'_{i})
$$
 
\end{enumerate}
\end{enumerate}
\end{definition}
 
By Lemma \ref{lemma-domain-stability}.\ref{lemma-domain-stability-01}, 
we have $\dom(R'_i) = \dom(R_i R) \subseteq \dom(R_i)$ for all $i \in [1,n]$. 
In the definition above, we have $i=n+1$ if and only if $\dom(R'_i) = \dom(R_i)$
for all $i \in [1,n]$. 

% 14:19 08/02/2026

\begin{theorem}[Transition Lemma]
\label{theorem-transition}
Suppose $\setR^+$ is injective, $\vec{R}$ is a sequence on $\setR^+$ and  
$\sigma \state \vec{R}$ and $R \in \setR^+$. Let $\tau = \trans(\sigma,R) $.
Then$\size(\sigma) > \size(\tau)$ and  $\tau \state \vec{R},R$.
\end{theorem}

\begin{proof}
Suppose $\setR^+$ is injective, $\vec{R}$ is a sequence on $\setR^+$ and
$\minStable(\vec{R}) = \vec{S} = S_1, \ldots, S_n$ be the minimum stable grouping
of $\vec{R}$. 
Let $S_{a_i}, \ldots, S_{b_i}$ for $i \in [1,h]$ be the components of $\vec{S}$. Assume
$\sigma = (U_1, d_1,V_1; \ldots; U_h, d_h, V_h) \state \vec{R}$: by definition of $\state$
we have $U_i = S_{a_i}, \ldots, S_{n}$ and $V_i = S_{b_i}$. 

Let $\tau = \trans(\sigma,R) $. Suppose
$R \in \setR^+$, in order to prove that $\tau \state \vec{R},R$. We argue by cases on the
definition of $\tau$. Let
$$
U'_i = U_i R = S_{a_i}, \ldots, S_{n} R \mbox{ for all $i \in [1,h]$ \ \ \ and } U'_{h+1} = R
$$ 
We have $\dom(U'_i) \subseteq \dom(U_i)$ 
by Lemma \ref{lemma-domain-inclusion-01}.\ref{lemma-domain-inclusion-01}.
Suppose that $i$ is the first element of $i \in [1,h+1]$ such that:

\begin{enumerate} 
\item
either $i \in [1,h]$ and $\dom(U'_i) \subset \dom(U_i)$
\item 
or $i=h+1$
\end{enumerate}

By definition unfolding, $i$ is the first index in $[1,h]$ such that 
$\dom(S_{a_i} \ldots S_{n} R) \subset \dom(S_{a_i} \ldots S_{n})$, if any, otherwise $i=h+1$.
By Lemma \ref{lemma-minstable-inductive}, 
in the case $i \in [1,h]$we have 
$$
\minStable(\vec{R},R) = 
S_{a_1}, \ldots, S_{b_1}, \ldots, S_{a_{i-1}}\ldots, S_{b_{i-1}}, (S_{a_i} \ldots S_{n} R)
$$,
in the case $i = h+1$ we have
$$
\minStable(\vec{R},R) = 
S_{a_1}, \ldots, S_{b_1}, \ldots, S_{a_{h}}\ldots, S_{b_{h}}, R
$$

We argue by cases on the definition of $\tau$.
\begin{enumerate}

\item 
\emph{Assume that $i=1$ or $\DomCard(U'_i) \not = \DomCard(U'_{i-1})$}and
$$
\trans(\sigma,R) = (U'_1, d_1, V_1; \ \ldots; U'_{i-1}, d_{i-1},V_{i-1}; \ U'_i, 0, \ U'_i)
$$

Suppose $i=1$. Then $\dom(U'_1) \subset \dom(U_1)$, therefore $\minStable(\vec{R},R) =
(S_{a_1} \ldots S_{n} R) = U'_1$, a grouping made of a single product, 
and $\trans(\sigma,R) = (U'_1, 0, U'_0)$.
$0$ is a degree of $U'_1$ w.r.t. $\setR^+$, therefore
$(U'_1, 0, U'_0)$ is a state of $\vec{R},R$.
From $\dom(U'_1) \subset \dom(U_1)$ we deduce $\size(U'_1) < \size(U_1)$,
hence $\sigma > \trans(\sigma,R)$. 
\\

Suppose $i \in [2,h]$. Then $\dom(U'_i) \subset \dom(U_i)$
 therefore $\minStable(\vec{R},R) = S_1, \ldots, S_{b_{i-1}}, (S_{a_i} \ldots S_n R)$
 and $\DomCard(U'_i) \not = \DomCard(U'_{i-1})$. We deduce that 
 $\DomCard(U'_i) > \DomCard(U'_{i-1})$, namely $\DomCard(S_{a_i} \ldots S_n R)
 > \DomCard(S_{a_{i-1}} \ldots S_{n} R) = $ 
 (by stability and 
 Lemma \ref{lemma-domain-stability}.\ref{lemma-domain-stability-02}) 
 $ \DomCard(S_{a_{i-1}} \ldots S_{b_{i-1}}$.
 Thus $(S_{b_i} \ldots S_n R)$
 is a component of $\vec{R},R$ made of a single product, therefore $\tau \state \vec{R},R$. 
From $\dom(U'_i) \subset \dom(U_i)$ we deduce $\size(U'_i) < \size(U_i)$,
hence $\sigma > \trans(\sigma,R)$.
\\

Suppose $i=h+1$. Then $\dom(U'_{h+1}) \subset \dom(U_{h+1})$
 therefore $\minStable(\vec{R},R) = S_1, \ldots, S_{b_{h}}, R$
 and $\DomCard(U'_{h+1}) \not = \DomCard(U'_{h})$. We deduce that 
 $\DomCard(R) = \DomCard(U'_{h+1}) > \DomCard(U'_{h}) =
  \DomCard(S_{a_{h}} \ldots S_{b_{h}} R) = $ 
 (by stability and 
 Lemma \ref{lemma-domain-stability}.\ref{lemma-domain-stability-02})  
 $ \DomCard(S_{a_{h}} \ldots S_{b_{h}}$. Namely 
 $R = U_{h+1}$ is a one-product component of $\vec{R},R$ and $\tau \state \vec{R},R$. 
We have $\sigma > \trans(\sigma,R)$ because $\trans(\sigma,R)$ is the
concatenation of $\sigma$ and of $(R,0,R)$.

\item
\emph{Assume that $i>1$ and $\DomCard(U'_i) = \DomCard(U'_{i-1})$}. 
\begin{enumerate}
\item
\emph{Assume $d_{i-1}\ge 1$ and
$V_{i-1}$ is progressing from $\dom(V_{i-1})$ to $\dom(U'_{i})$} and

$$
\trans(\sigma,S) = 
(U'_1, d_1,V_1;  \ \ldots; U'_{i-2}, d_{i-1},V_{i-2};  \ U'_{i-1}, 1, U'_{i-1})
$$

Suppose $i \in [2,h]$. By assumption, $\dom(U'_i) \subset \dom(U_i)$
 therefore $\minStable(\vec{R},R) = S_1, \ldots, S_{b_{i-1}}, (S_{a_i} \ldots S_n R)$
 and $\DomCard(U'_i) = \DomCard(U'_{i-1})$. We deduce that 
 $S_{a_{i-1}}, \ldots, S_{b_{i-1}}, (S_{a_i} \ldots S_n R)$
 is the last component of $\vec{R},R$ with last element $(S_{a_i} \ldots S_n R)$.
 $0$ is a degree of its product $U'_{i-1}$ in $\setR^+$.
 We deduce $\tau \state \vec{R},R$.
 $V_{i-1}$ is progressing from $\dom(V_{i-1})$ to $\dom(U'_{i})$. 
 By Lemma \ref{lemma-decreasing-size}we deduce that 
 $\size(U'_{i}) > \size(U'_{i}) \ge \size(U'_{i-1})$. 
 Thus, $\sigma > \tau$.
 \\
 
 Suppose $i=h+1$. By assumption, $\dom(R) = \dom(U'_{h+1}) \subset \dom(U_{h+1})$
 therefore $\minStable(\vec{R},R) = S_1, \ldots, S_{b_h}, R$
 and $\DomCard(U'_{h+1}) = \DomCard(U'_{h})$. We deduce that 
 $S_{a_h}, \ldots, S_{b_h}, R$
 is a component of $\vec{R},R$ with last element $R = U'_{h+1}$. 
 $0$ is a degree of its product  
 $U'_{h}$ in $\setR^+$. We deduce $\tau \state \vec{R},R$.
 $V_{h}$ is progressing from $\dom(V_{h})$ to $\dom(U'_{h+1})$
 By Lemma \ref{lemma-decreasing-size}we deduce that 
 $\size(U'_{h+1}) > \size(U'_{h+1}) \ge \size(U'_{h})$.  
 Thus, $\sigma > \tau$.
 
\item
\emph{Assume $d-{i-1}= 0$ or 
$V_{i-1}$ is stationary from $\dom(V_{i-1})$ to $\dom(U'_{i})$} and
$$
\trans(\sigma,S) = 
(U'_1, d_1,V_1;  \ \ldots; U'_{i-2}, d_{i-1},V_{i-2};  \ U'_{i-1}, d_{i-1}+1, U'_{i})
$$

If $i \in [2,h]$ then $S_{a_{i-1}}, \ldots, S_{b_{i-1}}, (S_{a_i} \ldots S_n R)$
 is the last component of $\vec{R},R$ with last element $(S_{a_i} \ldots S_n R)$. 
 By Lemma \ref{lemma-degree-constant},
  $d_{i-1}+1$ 
 is a degree of $U'_{i-1}$. We deduce $\tau \state \vec{R},R$.
 Let $c = \DomCard(U_{i-1}) = \DomCard(U'_{i-1}) $. We conclude that 
 $$
 \size(\sigma) = 
 (U_1, d_1, \ldots, U_{i-2}, d_{i-2}, U_{i-1}, \newton{\MaxInd}{c}-d_{i-1}) > 
  (U'_1, d_1, \ldots, U'_{i-2}, d_{i-2}, U'_{i-1}, \newton{\MaxInd}{c}-d_{i-1}-1) = 
  	size(\tau)
  	$$
 
If $i = h+1$ then $S_{a_{h}}, \ldots, S_{b_{h}}, R$ is the last component of $\vec{R},R$ with last element $R$. 
By Lemma \ref{lemma-degree-constant}, $d_{i-1}+1$ 
 is a degree of $U'_{i-1}$. We deduce $\tau \state \vec{R},R$.
 Let $c = \DomCard(U_{i-1}) = \DomCard(U'_{i-1}) $. We conclude that 
$$
 \size(\sigma) = 
(U_1, d_1, \ldots, U_{h-1}, d_{h-1},  U_{h},  \newton{\MaxInd}{c}-d_{i-1}) > 
(U'_1, d_1, \ldots, U'_{h-1}, d_{h-1}, U'_{h}, \newton{\MaxInd}{c}-d_{i-1}-1) = 
size(\tau)
$$
\end{enumerate}
\end{enumerate}
\end{proof}

We define a graph $\setG$ with transitions $\sigma \rightarrow_R \trans(\sigma,R)$ labeled
by $R \in \setR^+$.

%10:50 02/02/2026

The number of states is bounded by 
$
(2^{O(\MaxInd^2)} \times 2^\MaxInd \times 2^{O(\MaxInd^2)})
\times \ldots \times
(2^{O(\MaxInd^2)} \times 2^\MaxInd \times 2^{O(\MaxInd^2)})
$,
with at most $3\MaxInd$ product. 
We conclude that it is bounded by $2^{O(\MaxInd^3)}$. Each step requires time bounded by
$O(2^{\MaxInd^2}) \times \log_2(2^{O(\MaxInd^3)})$. 
We have $n \le 2^{\MaxInd^2}$, and cyclically 
for each $R_i$ with $i = 1, \ldots, n$ we add the state obtained by concatenating with 
$R_i$ and we compare it with the existing states sorted, in order to decide whether it is a new
state or not. We stop adding states in the first cycle we add no new state.
The total time cost is bounded by the number of relations times the cost of checking
whether we have a new state times the upper bound to the number of states.
The total is $O(2^{\MaxInd^2}) \times {O(\MaxInd^3)} \times 2^{O(\MaxInd^3)}$,
namely it is $2^{O(\MaxInd^3)}$ again. Therefore the total cost is poly-exponential.

%If $\vec{R} = R_1, \ldots, R_n \in \setR^+$ and $\vec{S}$ is the Minimum Stable
%grouping of $\vec{R} $ then we set $\size(\vec{R})  = \size(\vec{S}) $.

\section{A Ranking Function}
We define a ranking function. 


\subsection{An Exponential Time Translation From Cyclic To Inductive}
We define $Eq(\alpha, \beta)$ the set of equations $\size(x) > \size(y)$ such that the
the inductive formula number $x$ of $\alpha$ progresses to the inductive formula number $y$
 of $\beta$.
 
 We associate to $\alpha$ and $\vec{S}$ a list 
 $\trans(\alpha,\vec{S})=(l_1, \ldots, l_h)$ elements. If $c_i = \DomCard(R_{a_i})$,
 and $d$ is a stationary degree of $R_{a_i} \ldots R_{b_i}$
 then $l_i$ is
 $(
 \size(\alpha, R_{b_i}\ldots, R_n), 
 \newton{\MaxInd}{c_i} - d
 )$.
Alternatively, we can use a single integer index
$$
(\newton{\MaxInd}{c_i}+1)\size(\alpha, R_{b_i}\ldots, R_n) + (\newton{\MaxInd}{c_i} - d)
$$

We claim that if $\beta$ is a bud with companion $\alpha$, then 
$Eq(\alpha, \beta) \vdash \trans(\alpha,\vec{S}) > \trans(\beta,\vec{S})$

The minimum of all $\trans(\alpha,\vec{S})$ for all $\vec{S}$ is the required ranking function.
We represent $\vec{S}$ with $h$ pairs of the form $(R_{a_i}\ldots R_{b_i-1}),R_{b_i}$


\bibliographystyle{plain}
\bibliography{20241019-SCT-Cyclic}

 

 
\end{document}

 